% ============================================================================
% Memory Effects, Dissipation, and Current Density in TDDFT
% A Derivation-Based Graduate Tutorial
% ============================================================================
% Compile: pdflatex tddft_memory_tutorial && bibtex tddft_memory_tutorial &&
%          pdflatex tddft_memory_tutorial && pdflatex tddft_memory_tutorial
% ============================================================================
\documentclass[12pt,letterpaper]{article}

\usepackage{amsmath,amssymb,amsthm}
\usepackage{hyperref}
\usepackage{booktabs}
\usepackage{geometry}
\usepackage{enumitem}
\usepackage{tcolorbox}
\usepackage{bm}
\usepackage{xcolor}
\usepackage{mdframed}
\usepackage{mathtools}

\geometry{margin=1in}

\hypersetup{colorlinks=true,linkcolor=blue!70!black,citecolor=green!50!black}

% ---------- Boxed derivation environment ----------
\tcbuselibrary{skins,breakable}
\newtcolorbox{derivbox}[1][]{
  colback=blue!3!white, colframe=blue!40!black,
  fonttitle=\bfseries\small, title={Derivation},
  breakable, #1
}
\newtcolorbox{keyresult}[1][]{
  colback=green!4!white, colframe=green!50!black,
  fonttitle=\bfseries\small, title={Key Result},
  breakable, #1
}
\newtcolorbox{insight}[1][]{
  colback=orange!5!white, colframe=orange!60!black,
  fonttitle=\bfseries\small, title={Physical Insight},
  breakable, #1
}

% ---------- Macros ----------
\newcommand{\bra}[1]{\langle #1 \rvert}
\newcommand{\ket}[1]{\lvert #1 \rangle}
\newcommand{\braket}[2]{\langle #1 \vert #2 \rangle}
\newcommand{\mel}[3]{\langle #1 \vert #2 \vert #3 \rangle}
\newcommand{\comm}[2]{\left[#1,\,#2\right]}
\newcommand{\dd}{\mathrm{d}}
\newcommand{\ii}{\mathrm{i}}
\newcommand{\Hks}{H_{\mathrm{KS}}}
\newcommand{\Vxc}{V_{\mathrm{xc}}}
\newcommand{\fxc}{f_{\mathrm{xc}}}
\newcommand{\Exc}{E_{\mathrm{xc}}}
\newcommand{\etal}{\textit{et al.}}

\theoremstyle{definition}
\newtheorem{theorem}{Theorem}
\newtheorem{remark}{Remark}

% ============================================================================
\begin{document}

\title{\Large\bfseries Memory Effects, Dissipation, and Current Density\\
in Time-Dependent Density Functional Theory\\[6pt]
\normalsize\textit{A Derivation-Based Graduate Tutorial}}

\author{Notes compiled from discussion --- Jacek Jakowski}
\date{\today}
\maketitle

\begin{abstract}
This tutorial addresses four questions that arise naturally when studying
RT-TDDFT from a quantum-chemistry perspective:
(i) what does the adiabatic approximation discard and why finite-order
time-derivative corrections are insufficient;
(ii) why a closed unitary quantum system can still exhibit dephasing in
the one-body density matrix;
(iii) why double excitations are absent from adiabatic TDDFT and what
is actually required to recover them;
(iv) why the charge density alone is insufficient as a basic variable
for periodic systems and how current-density TDDFT resolves this.
Each topic is developed through explicit derivations at the level of a
first-year graduate student in quantum chemistry or condensed-matter physics.
\end{abstract}

\tableofcontents
\newpage

% ============================================================================
\section{Background: The Exact TDDFT Framework}
\label{sec:background}
% ============================================================================

We work in atomic units ($\hbar = m_e = e = 1$) throughout.
The time-dependent Kohn-Sham (TDKS) equations are:
\begin{equation}
  \ii\frac{\partial\psi_n}{\partial t} = \Hks[\rho]\,\psi_n,
  \qquad
  \Hks = -\tfrac{1}{2}\nabla^2 + V_{\mathrm{ext}} + V_H[\rho] + \Vxc[\rho],
  \label{eq:tdks}
\end{equation}
with electron density $\rho(\bm{r},t) = \sum_n f_n|\psi_n|^2$.
The \textbf{occupation numbers} $f_n \in \{0,1\}$ (or $\{0,1,2\}$ for
spin-restricted closed-shell) are set by the \textbf{Aufbau principle}
at $t=0$: orbitals are filled from lowest to highest KS eigenvalue
$\epsilon_n$, giving $f_n = 1$ for the $N$ lowest and $f_n = 0$ for
the rest.
They are \emph{fixed at their initial values} for all $t>0$: the TDKS
equations propagate the orbital shapes $\psi_n(\bm{r},t)$ but do not
change which orbitals are occupied.
This is why adiabatic TDDFT cannot describe $T_1$ population relaxation
(Section~\ref{sec:dissipation}) --- the $f_n$ are frozen constants, not
dynamical variables.
(Koopmans' theorem connects $\epsilon_n$ to ionisation energies;
Janak's theorem generalises this to fractional occupations, relevant for
smearing in metallic systems, but not used in the present tutorial.)

\paragraph{The Runge-Gross theorem.}
Recall the \textbf{Hohenberg-Kohn (HK) theorem} (1964): for a
ground-state system, there is a one-to-one mapping between the static
external potential $V_{\mathrm{ext}}(\bm{r})$ and the ground-state density
$\rho_0(\bm{r})$.
Consequently, every ground-state observable is a functional of $\rho_0$
alone — this is the foundation of ground-state DFT.

Runge and Gross (1984) proved the time-dependent analogue of HK: given a fixed
initial state $\Psi_0$, there is
a one-to-one mapping between the time-dependent external potential
$V_{\mathrm{ext}}(\bm{r},t)$ and the density $\rho(\bm{r},t)$.
This guarantees the existence of an exact time-dependent XC functional:
\begin{equation}
  \Vxc(\bm{r},t) = \Vxc\!\left[\rho(\bm{r}',t');\,t'\leq t\right](\bm{r},t).
  \label{eq:exact_vxc}
\end{equation}
The critical feature of Eq.~(\ref{eq:exact_vxc}) is that $\Vxc$ at time $t$
is a \emph{functional of the entire density history} from the initial time
to $t$ --- not just a function of the instantaneous density $\rho(\bm{r},t)$.

\paragraph{The adiabatic approximation defined.}
The adiabatic approximation replaces the history-dependent functional by
the ground-state XC functional evaluated at the instantaneous density:
\begin{equation}
  \boxed{\Vxc^{\mathrm{ATDDFT}}(\bm{r},t) = \Vxc^{\mathrm{gs}}\!\left[\rho(\cdot,t)\right](\bm{r})}.
  \label{eq:atddft}
\end{equation}
At every timestep, Eq.~(\ref{eq:atddft}) calls the same LDA/GGA routine
used in a ground-state calculation, passing the instantaneous density.
The functional has no knowledge that $\rho$ is time-evolving.

% ============================================================================
\section{Functional Derivatives: A Primer for Quantum Chemists}
\label{sec:funcderiv}
% ============================================================================

This section is a self-contained introduction to functional derivatives,
written for readers comfortable with Gaussian-basis DFT but not variational
calculus. Every symbol used later in this tutorial---$\Vxc = \delta\Exc/\delta\rho$,
the XC kernel $\fxc = \delta\Vxc/\delta\rho$, the linear-response equations---is
a functional derivative. Understanding what these objects \emph{are} makes the
rest of the theory transparent.

\begin{tcolorbox}[colback=yellow!6!white, colframe=orange!70!black,
  fonttitle=\bfseries\small, title={Notation: Three Different Uses of $\delta$},
  breakable]
The symbol $\delta$ carries three distinct meanings in this section.
They are easy to tell apart by context, but it is worth listing them
explicitly before they all appear on the same page.

\medskip
\begin{tabular}{@{}lll@{}}
\toprule
Symbol & Name & What it is \\
\midrule
$\dfrac{\delta F}{\delta f(x)}$ & functional derivative
  & a \emph{function} of $x$; measures sensitivity of $F$ to $f$ at $x$ \\[10pt]
$\delta(x - x_0)$ & Dirac delta
  & a \emph{distribution}; zero everywhere except $x_0$, unit integral \\[10pt]
$\delta f$, $\delta\rho$ & variation
  & an \emph{infinitesimal perturbation} of the function $f$ or density $\rho$ \\
\bottomrule
\end{tabular}

\medskip
\noindent The functional derivative $\delta F/\delta f(x)$ is defined via an
integral against a test function $\phi$
(Eq.~\ref{eq:fd_def} below): it is never a ratio of two $\delta$'s.
The Dirac delta $\delta(x-x_0)$ appears when that test function is chosen
to be a spike at a single point (Eq.~\ref{eq:fd_pointwise}).
A variation $\delta f$ is the perturbation you add to $f$; it plays
the role of $\phi$ when no specific form is needed.
\end{tcolorbox}

% ----------------------------------------------------------------------------
\subsection{What Is a Functional?}
% ----------------------------------------------------------------------------

In ordinary calculus a \emph{function} maps a number to a number:
$f : \mathbb{R} \to \mathbb{R}$, e.g.\ $f(x) = x^2$.
A \emph{functional} maps an entire \emph{function} to a number:
$F : \{\text{functions}\} \to \mathbb{R}$.

\begin{center}
\begin{tabular}{lll}
\toprule
Object & Input & Output \\
\midrule
Function $f(x)$ & a number $x$ & a number \\
Functional $F[f]$ & an entire function $f(\cdot)$ & a number \\
\bottomrule
\end{tabular}
\end{center}

\noindent The square bracket $F[\,\cdot\,]$ is the conventional notation for a
functional to distinguish it from a function $f(\,\cdot\,)$.

\paragraph{Examples you already know.}
\begin{itemize}[leftmargin=*]
  \item \textbf{Norm}: $\|f\|^2 = \int |f(x)|^2\,\dd x$ --- maps a wavefunction to a
    real number.
  \item \textbf{Energy expectation value}: $E[\Psi] = \langle\Psi|H|\Psi\rangle$ ---
    maps a many-body wavefunction to the energy.
  \item \textbf{Hartree energy}: $E_H[\rho] = \tfrac{1}{2}\iint
    \frac{\rho(\bm{r})\,\rho(\bm{r}')}{|\bm{r}-\bm{r}'|}\,\dd^3r\,\dd^3r'$ ---
    maps the electron density to the classical electrostatic self-energy.
  \item \textbf{XC energy}: $\Exc[\rho]$ --- maps $\rho$ to a number whose exact
    form is unknown; approximated by LDA, PBE, etc.
\end{itemize}

% ----------------------------------------------------------------------------
\subsection{The Functional Derivative: Intuition First}
% ----------------------------------------------------------------------------

\paragraph{Analogy with ordinary derivatives.}
For a function $f(x)$, the derivative $f'(x_0)$ answers:
\begin{quote}
  \textit{If I nudge the argument by $\varepsilon$ at the point $x_0$,
  how much does $f$ change?}
\end{quote}
$\delta f = f'(x_0)\,\delta x$.

For a functional $F[f]$, the functional derivative
$\frac{\delta F}{\delta f(x_0)}$ answers:
\begin{quote}
  \textit{If I nudge the input function $f$ by a tiny spike at the
  single point $x_0$, how much does $F$ change?}
\end{quote}

\begin{keyresult}[title={The Dirac Delta: Definition and Sifting Property}]
The \textbf{Dirac delta} $\delta(x - x_0)$ is not a function in the ordinary
sense but a \emph{distribution} (generalised function) defined by what it does
inside an integral.  Its two essential properties are:
\begin{align}
  \delta(x - x_0) &= 0 \quad \text{for all } x \neq x_0,
  \label{eq:dirac_zero}\\
  \int_{-\infty}^{\infty} \delta(x - x_0)\,g(x)\,\dd x &= g(x_0)
  \quad\text{for any continuous } g.
  \label{eq:sifting}
\end{align}
Equation~(\ref{eq:sifting}) is the \textbf{sifting property}: the delta
``picks out'' the value of $g$ at the single point $x_0$.

\medskip
\noindent\textbf{Contrast with the Kronecker delta.}
The Kronecker delta $\delta_{ij}$ is an ordinary number (0 or 1) that acts on
discrete indices.  The Dirac delta $\delta(x-x_0)$ is a continuous analogue:
it is zero everywhere except at one point, but it is normalised so that its
integral equals 1.  On a grid with spacing $h$, $\delta(x-x_0)\approx \delta_{i,i_0}/h$,
which recovers the Kronecker delta in the discrete limit.

\medskip
\noindent\textbf{Intuitive picture.}
Think of $\delta(x-x_0)$ as the limit of a very narrow, very tall Gaussian:
$\delta(x-x_0) = \lim_{\sigma\to 0}
\frac{1}{\sigma\sqrt{2\pi}}\exp\!\bigl[-(x-x_0)^2/2\sigma^2\bigr]$.
The area under the spike is always 1, but the width shrinks to zero.
\end{keyresult}

The ``tiny spike at $x_0$'' is mathematically a Dirac delta:
$f(x) \to f(x) + \varepsilon\,\delta(x - x_0)$.
This is the infinitesimal perturbation of a function: it leaves $f$
unchanged everywhere except at $x_0$, where it adds $\varepsilon$.

\paragraph{Gradient analogy.}
If $F$ depended on a finite set of numbers $(a_1, a_2, \ldots, a_N)$
instead of on a continuous function $f(x)$, its sensitivity to each
$a_i$ would be the partial derivative $\partial F/\partial a_i$, and
the full gradient would be
$\nabla F = (\partial F/\partial a_1, \ldots, \partial F/\partial a_N)$.

A functional $F[f]$ is the limit $N\to\infty$ of this construction,
where $a_i \equiv f(x_i)$ are the values of $f$ at a dense grid of
points. The functional derivative $\delta F/\delta f(x)$ is the
continuum analogue of the gradient component at grid point $x$:

\begin{equation}
  \frac{\delta F}{\delta f(x)} \;\longleftrightarrow\;
  \frac{\partial F}{\partial a_i}\bigg|_{a_i = f(x_i)},
  \qquad x_i \to x.
  \label{eq:fd_gradient_analogy}
\end{equation}

This is not merely suggestive --- it is exact on a discrete grid.
\begin{derivbox}[title={Proof: functional derivative = partial derivative / grid spacing}]
Let the grid have $N$ points $x_1,\ldots,x_N$ with uniform spacing $h$,
so $f$ is represented by the vector $(a_1,\ldots,a_N)$ with $a_i = f(x_i)$.
A functional $F[f] = \int g(f(x))\,\dd x$ is approximated by
$F \approx \sum_k g(a_k)\,h$.
Apply the definition~(\ref{eq:fd_def}) with $\phi(x) = \delta(x-x_i)$
(a spike at grid point $x_i$, which on the grid acts as a unit vector in
the $i$-th direction):
\begin{align}
  \frac{\delta F}{\delta f(x_i)}
  &= \lim_{\varepsilon\to 0}
     \frac{F[\bm{a}+\varepsilon\bm{e}_i] - F[\bm{a}]}{\varepsilon}
   = \frac{\partial F}{\partial a_i}
   = h\,g'(a_i).
\end{align}
But the definition also requires dividing by the ``volume element'' to
obtain a density: comparing $\int(\delta F/\delta f(x))\phi(x)\,\dd x$
with the finite sum $\sum_k (\partial F/\partial a_k)\phi_k$ shows
that $\delta F/\delta f(x_i) = (\partial F/\partial a_i)/h$.

For the example $F=\sum_k a_k^2 h$:
$\partial F/\partial a_i = 2a_i h$, so
$\delta F/\delta f(x_i) = 2a_i h / h = 2a_i = 2f(x_i)$.
This is exactly the continuum result $\delta F/\delta f(x)=2f(x)$,
confirmed at grid level.
\end{derivbox}

% ----------------------------------------------------------------------------
\subsection{Formal Definition}
% ----------------------------------------------------------------------------

\begin{keyresult}[title={Definition: Functional Derivative}]
The functional derivative of $F[f]$ with respect to $f$ at the point
$x_0$ is defined by:
\begin{equation}
  \int \frac{\delta F}{\delta f(x)}\,\phi(x)\,\dd x
  \;=\;
  \lim_{\varepsilon\to 0}
  \frac{F[f + \varepsilon\phi] - F[f]}{\varepsilon}
  \;\equiv\;
  \left.\frac{\dd}{\dd\varepsilon}F[f+\varepsilon\phi]\right|_{\varepsilon=0},
  \label{eq:fd_def}
\end{equation}
for all ``\textbf{test functions}'' $\phi$.
A test function is any well-behaved function used to \emph{probe} the
sensitivity of $F$ to changes in $f$: it represents the direction in
function space along which $f$ is perturbed.
Any choice of $\phi$ gives a valid directional derivative; different
choices probe different directions.
In a Gaussian-basis code, the discrete analogue of $\phi$ is a variation
$\delta P_{\mu\nu}$ of one density-matrix element: perturbing $P_{\mu\nu}$
by $\varepsilon$ probes the sensitivity of $E[\bm{P}]$ to that element,
giving $\partial E/\partial P_{\mu\nu}$ --- exactly the matrix-element
version of the functional derivative.
Choosing $\phi(x) = \delta(x-x_0)$
and applying the sifting property~(\ref{eq:sifting}) to the left-hand side:
\begin{equation}
  \int \frac{\delta F}{\delta f(x)}\,\delta(x-x_0)\,\dd x
  = \frac{\delta F}{\delta f(x_0)},
\end{equation}
recovers the pointwise version:
\begin{equation}
  \frac{\delta F}{\delta f(x_0)}
  = \left.\frac{\dd}{\dd\varepsilon}F\bigl[f+\varepsilon\,\delta(\cdot-x_0)\bigr]
    \right|_{\varepsilon=0}.
  \label{eq:fd_pointwise}
\end{equation}
\end{keyresult}

\begin{keyresult}[title={Algorithm: How to Compute a Functional Derivative}]
Every functional derivative calculation in this tutorial follows the same
four steps.
\begin{enumerate}[label=\textbf{Step \arabic*.}, leftmargin=*, itemsep=4pt]
  \item \textbf{Perturb.}
    Replace $f \to f + \varepsilon\phi$ in the definition of $F$ and expand
    to first order in $\varepsilon$, discarding $O(\varepsilon^2)$ terms.
  \item \textbf{Subtract and divide.}
    Form $(F[f+\varepsilon\phi] - F[f])/\varepsilon$ and take $\varepsilon\to 0$
    (equivalently, $\dd/\dd\varepsilon|_{\varepsilon=0}$).
    The result is a linear functional of $\phi$.
  \item \textbf{Write as a single integral.}
    Bring the result into the form $\int (\cdots)\,\phi(x)\,\dd x$,
    using symmetry of the integration variable or Fubini's theorem
    (switching order of a double integral) if needed.
  \item \textbf{Read off the coefficient.}
    The expression multiplying $\phi(x)$ inside the integral is
    $\delta F/\delta f(x)$.
    That is the functional derivative.
\end{enumerate}
\noindent This procedure defines the \textbf{Gâteaux derivative} (named after
René Gâteaux, 1913) --- the functional analogue of the directional derivative
in ordinary calculus.
It is a weaker notion than the Fréchet derivative (which requires uniform
convergence over all directions), but for the smooth, well-behaved functionals
that arise in DFT the two coincide.
No topology or functional-analysis background is required; the four-step
algorithm above is all one ever needs in practice.
\end{keyresult}

% ----------------------------------------------------------------------------
\subsection{Worked Example 1 --- A Simple Integral Functional}
% ----------------------------------------------------------------------------

Let
\begin{equation}
  F[f] = \int f(x)^2\,\dd x.
  \label{eq:ex1_F}
\end{equation}

\begin{derivbox}
Perturb $f \to f + \varepsilon\phi$:
\begin{align}
  F[f+\varepsilon\phi]
  &= \int \bigl(f(x)+\varepsilon\phi(x)\bigr)^2\,\dd x \notag\\
  &= \int f(x)^2\,\dd x
     + 2\varepsilon\int f(x)\,\phi(x)\,\dd x
     + \varepsilon^2\int\phi(x)^2\,\dd x.
\end{align}
Subtract $F[f]$, divide by $\varepsilon$, and take $\varepsilon\to 0$:
\begin{equation}
  \left.\frac{\dd}{\dd\varepsilon}F[f+\varepsilon\phi]\right|_{\varepsilon=0}
  = 2\int f(x)\,\phi(x)\,\dd x.
\end{equation}
Comparing with the definition (\ref{eq:fd_def}),
the coefficient of $\phi(x)$ in the integrand is $2f(x)$, so:
\begin{equation}
  \boxed{\frac{\delta F}{\delta f(x)} = 2f(x).}
\end{equation}
\end{derivbox}

\noindent\textbf{Check with ordinary calculus.}
On a finite grid of $N$ points with spacing $h$,
$F \approx \sum_i f_i^2\,h$, so $\partial F/\partial f_j = 2f_j\,h$.
Dividing by $h$ to get the continuum limit gives $2f(x_j)$. Consistent.

% ----------------------------------------------------------------------------
\subsection{Worked Example 2 --- The Hartree Energy}
% ----------------------------------------------------------------------------

The Hartree (Coulomb self-interaction) energy is:
\begin{equation}
  E_H[\rho] = \frac{1}{2}\iint
    \frac{\rho(\bm{r})\,\rho(\bm{r}')}{|\bm{r}-\bm{r}'|}
    \,\dd^3r\,\dd^3r'.
  \label{eq:hartree}
\end{equation}
The factor of $\tfrac{1}{2}$ prevents \textbf{double counting}.
The integrand $\rho(\bm{r})\rho(\bm{r}')/|\bm{r}-\bm{r}'|$ counts the
interaction between volume elements at $\bm{r}$ and $\bm{r}'$.
Because the double integral $\int\!\int\!\dd^3r\,\dd^3r'$ runs over
all ordered pairs, each unordered pair $\{\bm{r},\bm{r}'\}$ appears
twice (once as $(\bm{r},\bm{r}')$ and once as $(\bm{r}',\bm{r})$).
The $\tfrac{1}{2}$ divides by 2 to correct this.
This is the \emph{nonlocal} functional: $E_H$ depends on $\rho$ at
\emph{all} pairs of points, not just locally.

\begin{derivbox}
Perturb $\rho \to \rho + \varepsilon\phi$ and expand the product:
\begin{align}
  E_H[\rho+\varepsilon\phi]
  &= \frac{1}{2}\iint
    \frac{(\rho(\bm{r})+\varepsilon\phi(\bm{r}))
          (\rho(\bm{r}')+\varepsilon\phi(\bm{r}'))}
         {|\bm{r}-\bm{r}'|}
    \,\dd^3r\,\dd^3r' \notag\\
  &= E_H[\rho]
     + \frac{\varepsilon}{2}\iint
       \frac{\rho(\bm{r}')\,\phi(\bm{r})}{|\bm{r}-\bm{r}'|}
       \,\dd^3r\,\dd^3r'
     + \frac{\varepsilon}{2}\iint
       \frac{\rho(\bm{r})\,\phi(\bm{r}')}{|\bm{r}-\bm{r}'|}
       \,\dd^3r\,\dd^3r'
     + O(\varepsilon^2).
\end{align}
\textbf{Relabelling dummy variables.}
In the second $\varepsilon$-linear term, swap the integration variables
$\bm{r}\leftrightarrow\bm{r}'$ (both are dummy variables of integration;
$|\bm{r}-\bm{r}'|=|\bm{r}'-\bm{r}|$):
\begin{equation}
  \frac{1}{2}\iint\frac{\rho(\bm{r})\,\phi(\bm{r}')}{|\bm{r}-\bm{r}'|}
  \,\dd^3r\,\dd^3r'
  \;\overset{\bm{r}\leftrightarrow\bm{r}'}{=}\;
  \frac{1}{2}\iint\frac{\rho(\bm{r}')\,\phi(\bm{r})}{|\bm{r}-\bm{r}'|}
  \,\dd^3r\,\dd^3r'.
\end{equation}
Both $\varepsilon$-linear terms are now identical; their sum gives a factor
of 2 that cancels the $\tfrac{1}{2}$:
\begin{equation}
  E_H[\rho+\varepsilon\phi]
  = E_H[\rho]
    + \varepsilon\iint
      \frac{\rho(\bm{r}')\,\phi(\bm{r})}{|\bm{r}-\bm{r}'|}
      \,\dd^3r\,\dd^3r'
    + O(\varepsilon^2).
\end{equation}
\textbf{Fubini's theorem} (switching the order of the $\bm{r}$ and $\bm{r}'$
integrations, valid here because the integrand is absolutely integrable away
from the diagonal) lets us write the double integral as an iterated integral
with the $\bm{r}'$ integration performed first:
\begin{equation}
  \left.\frac{\dd}{\dd\varepsilon}E_H[\rho+\varepsilon\phi]\right|_{\varepsilon=0}
  = \int \phi(\bm{r})
    \underbrace{\left[\int\frac{\rho(\bm{r}')}{|\bm{r}-\bm{r}'|}\,\dd^3r'\right]}_{V_H(\bm{r})}
    \,\dd^3r.
\end{equation}
The term in brackets is exactly the Hartree potential $V_H(\bm{r})$. So:
\begin{equation}
  \boxed{\frac{\delta E_H}{\delta\rho(\bm{r})} = V_H(\bm{r})
    = \int\frac{\rho(\bm{r}')}{|\bm{r}-\bm{r}'|}\,\dd^3r'.}
  \label{eq:vh_fd}
\end{equation}
\end{derivbox}

\noindent\textbf{Physical meaning.}
$V_H(\bm{r})$ is the electrostatic potential at $\bm{r}$ due to the
charge distribution $\rho$. The functional derivative says: if you add
one electron at $\bm{r}_0$, the electrostatic energy increases by $V_H(\bm{r}_0)$.
This is exactly the definition of electrostatic potential from classical
physics.

\begin{insight}
\textbf{Connection to Gaussian-basis DFT: from $\delta E_H/\delta\rho$ to
$\partial E_H/\partial P_{\mu\nu}$.}
In a Gaussian-basis code the electron density is written
$\rho(\bm{r}) = \sum_{\mu\nu}P_{\mu\nu}\,\chi_\mu(\bm{r})\chi_\nu(\bm{r})$,
where $P_{\mu\nu}$ is the one-particle density matrix and $\{\chi_\mu\}$ are
the basis functions.
The Hartree energy becomes a function of the matrix elements $P_{\mu\nu}$:
\begin{equation}
  E_H = \tfrac{1}{2}\sum_{\mu\nu\lambda\sigma}
        P_{\mu\nu}\,P_{\lambda\sigma}\,(\mu\nu|\lambda\sigma),
  \label{eq:hartree_gaussian}
\end{equation}
where $(\mu\nu|\lambda\sigma)=\iint
\chi_\mu(\bm{r})\chi_\nu(\bm{r})\,|\bm{r}-\bm{r}'|^{-1}\,
\chi_\lambda(\bm{r}')\chi_\sigma(\bm{r}')\,\dd^3r\,\dd^3r'$
is the two-electron repulsion integral (ERI), familiar from Hartree-Fock.

The analogue of $\delta E_H/\delta\rho = V_H$ is the \emph{Coulomb matrix}
$\bm{G}$, obtained by differentiating with respect to the density matrix:
\begin{equation}
  G_{\mu\nu} \equiv \frac{\partial E_H}{\partial P_{\mu\nu}}
  = \sum_{\lambda\sigma} P_{\lambda\sigma}\,(\mu\nu|\lambda\sigma).
  \label{eq:G_matrix}
\end{equation}
This is a matrix equation ($\bm{G} = \bm{J}[\bm{P}]$ in standard notation).
The continuous functional derivative $\delta E_H/\delta\rho(\bm{r}) = V_H(\bm{r})$
and the discrete matrix derivative $\partial E_H/\partial P_{\mu\nu} = G_{\mu\nu}$
are the same object in two different representations:
$V_H(\bm{r}) = \sum_{\mu\nu}G_{\mu\nu}\chi_\mu(\bm{r})\chi_\nu(\bm{r})$.

\medskip
\noindent The real-space grid in RMG plays exactly the role of the Gaussian
basis: the continuous functional derivative $\delta E_H/\delta\rho(\bm{r}_i)$
at grid point $\bm{r}_i$ is the grid-point analogue of $G_{\mu\nu}$.
\end{insight}

\begin{insight}
The Kohn-Sham equation is derived by requiring that the total energy
$E[\rho]$ is stationary under variations $\delta\rho$ subject to
the constraint $\int\rho\,\dd^3r = N$.
Setting $\delta E/\delta\rho = \mu$ (the chemical potential) gives
the Euler equation, from which the KS potential emerges as
$V_\text{KS} = V_\text{ext} + \delta E_H/\delta\rho + \delta\Exc/\delta\rho$.
Functional derivatives are not a formalism bolted onto DFT --- they
\emph{define} the KS potential.
\end{insight}

% ----------------------------------------------------------------------------
\subsection{Worked Example 3 --- The Bilinear Functional}
% ----------------------------------------------------------------------------

The Hartree energy is a special case of the \emph{bilinear functional}:
\begin{equation}
  F[f] = \iint K(x,y)\,f(x)\,f(y)\,\dd x\,\dd y,
  \label{eq:bilinear_F}
\end{equation}
where $K(x,y)$ is a fixed symmetric kernel, $K(x,y)=K(y,x)$.
(For Hartree: $f=\rho$, $K(\bm{r},\bm{r}')=\tfrac{1}{2}|\bm{r}-\bm{r}'|^{-1}$.)
This example is worth working through because its first and second functional
derivatives appear in three places later in this tutorial: the Hartree potential,
the Casida coupling matrix, and hybrid XC kernels.

\paragraph{First functional derivative.}
\begin{derivbox}
Perturb $f \to f + \varepsilon\phi$:
\begin{align}
  F[f+\varepsilon\phi]
  &= \iint K(x,y)\bigl(f(x)+\varepsilon\phi(x)\bigr)
                    \bigl(f(y)+\varepsilon\phi(y)\bigr)\,\dd x\,\dd y \notag\\
  &= F[f]
     + \varepsilon\iint K(x,y)\,\phi(x)\,f(y)\,\dd x\,\dd y
     + \varepsilon\iint K(x,y)\,f(x)\,\phi(y)\,\dd x\,\dd y
     + O(\varepsilon^2).
\end{align}
Rename $x\leftrightarrow y$ in the second integral and use $K(y,x)=K(x,y)$:
both integrals are equal.  Taking $\varepsilon\to 0$:
\begin{equation}
  \left.\frac{\dd}{\dd\varepsilon}F[f+\varepsilon\phi]\right|_{\varepsilon=0}
  = 2\int\phi(x)\underbrace{\left[\int K(x,y)\,f(y)\,\dd y\right]}_{
    \equiv\,K*f\,(x)}\dd x.
\end{equation}
Reading off the coefficient of $\phi(x)$:
\begin{equation}
  \boxed{\frac{\delta F}{\delta f(x)} = 2\int K(x,y)\,f(y)\,\dd y.}
  \label{eq:bilinear_first}
\end{equation}
\end{derivbox}

\noindent\textbf{Check against Hartree.}
With $K(\bm{r},\bm{r}')=\tfrac{1}{2}|\bm{r}-\bm{r}'|^{-1}$, the factor of 2
cancels the $\tfrac{1}{2}$ and Eq.~(\ref{eq:bilinear_first}) reduces to
$\delta E_H/\delta\rho(\bm{r}) = \int \rho(\bm{r}')/|\bm{r}-\bm{r}'|\,\dd^3r' = V_H(\bm{r})$,
consistent with Eq.~(\ref{eq:vh_fd}).

\paragraph{Second functional derivative.}
Apply Eq.~(\ref{eq:sfd_def}) to $\delta F/\delta f(x) = 2\int K(x,y)f(y)\,\dd y$.
This is a linear functional of $f$, so the second derivative is simply its
kernel:
\begin{derivbox}
Perturb $f\to f+\varepsilon\,\delta(\cdot-x')$ inside
$\delta F/\delta f(x) = 2\int K(x,y)f(y)\,\dd y$:
\begin{align}
  \frac{\delta}{\delta f(x')}\left[2\int K(x,y)f(y)\,\dd y\right]
  &= 2\int K(x,y)\,\delta(y-x')\,\dd y \notag\\
  &= 2K(x,x').
\end{align}
Therefore:
\begin{equation}
  \boxed{\frac{\delta^2 F}{\delta f(x)\,\delta f(x')} = 2K(x,x').}
  \label{eq:bilinear_second}
\end{equation}
\end{derivbox}

\noindent\textbf{Interpretation.}
The second functional derivative of a bilinear functional is (twice) the kernel
itself.  The kernel $K(x,x')$ directly measures the nonlocality of $F$:
\begin{itemize}[leftmargin=*]
  \item If $K(x,y)=c\,\delta(x-y)$ (local kernel), $\delta^2F/\delta f^2 = 2c\,\delta(x-x')$
    --- a diagonal Hessian, like LDA.
  \item If $K(x,y)=1/|x-y|$ (Coulomb), $\delta^2 F/\delta f^2 = 2/|x-x'|$
    --- fully nonlocal, like exact exchange.
\end{itemize}

\begin{insight}
\textbf{Connection to Casida and hybrid XC.}
In the Casida linear-response equations (Section~\ref{sec:casida}),
the coupling matrix element is
$K_{ia,jb} = \iint\phi_i(\bm{r})\phi_a(\bm{r})
\left[\frac{1}{|\bm{r}-\bm{r}'|}+\fxc(\bm{r},\bm{r}')\right]
\phi_j(\bm{r}')\phi_b(\bm{r}')\,\dd^3r\,\dd^3r'$.
The Coulomb part is structurally $2K(x,x')$ from Eq.~(\ref{eq:bilinear_second})
with $K = 1/|\bm{r}-\bm{r}'|$ and the density evaluated in MO pairs.
Adding $\fxc$ is the XC correction to this bilinear structure.
In hybrid functionals, a fraction $\alpha$ of the exact exchange contributes
$2\alpha K(\bm{r},\bm{r}')$ with $K=1/2|\bm{r}-\bm{r}'|$ --- again bilinear.
\end{insight}

% ----------------------------------------------------------------------------
\subsection{The XC Potential as a Functional Derivative}
% ----------------------------------------------------------------------------

The total ground-state energy in DFT is decomposed as:
\begin{equation}
  E_{\mathrm{tot}}[\rho]
  = T_s[\rho] + E_{\mathrm{ext}}[\rho] + E_H[\rho] + \Exc[\rho],
  \label{eq:etot_decomp}
\end{equation}
where $T_s[\rho]$ is the kinetic energy of the \emph{non-interacting}
KS electrons (not the true interacting kinetic energy),
$E_{\mathrm{ext}}[\rho] = \int V_{\mathrm{ext}}(\bm{r})\rho(\bm{r})\,\dd^3r$
is the electron-nuclear attraction, and $E_H[\rho]$ is the Hartree
energy~(\ref{eq:hartree}).

The XC energy $\Exc[\rho]$ is defined as whatever is left over:
\begin{equation}
  \Exc[\rho]
  \equiv E_{\mathrm{tot}}^{\mathrm{exact}}[\rho]
       - T_s[\rho] - E_{\mathrm{ext}}[\rho] - E_H[\rho].
  \label{eq:exc_def}
\end{equation}
It collects the kinetic correlation (difference between true and
non-interacting kinetic energies) and all exchange-correlation effects
beyond the classical Hartree repulsion.
Its explicit form is unknown, but its functional derivative defines the XC
potential that appears in the KS equation:

\begin{equation}
  \boxed{\Vxc(\bm{r}) = \frac{\delta\Exc[\rho]}{\delta\rho(\bm{r})}.}
  \label{eq:vxc_fd}
\end{equation}

\paragraph{LDA: functional derivative reduces to an ordinary derivative.}
In the local density approximation (LDA):
\begin{equation}
  \Exc^{\mathrm{LDA}}[\rho] = \int \varepsilon_{\mathrm{xc}}\!\left(\rho(\bm{r})\right)
    \rho(\bm{r})\,\dd^3r,
  \label{eq:exc_lda}
\end{equation}
where $\varepsilon_{\mathrm{xc}}(\rho)$ is the XC energy per particle of a
\emph{uniform} electron gas at density $\rho$.
The density of the uniform gas is conveniently parametrised by the
\textbf{Wigner-Seitz radius} $r_s = (3/4\pi\rho)^{1/3}$ (in bohr),
the radius of the sphere containing one electron on average.
For real materials $r_s \approx 2$--$6\,a_0$ (metals: $r_s\approx2$--$3$;
valence electrons in atoms: $r_s\approx 3$--$6$).
The exact $\varepsilon_{\mathrm{xc}}(r_s)$ is known from
\textbf{quantum Monte Carlo} (Ceperley and Alder, 1980) and fitted to
analytic forms: the \textbf{VWN} (Vosko-Wilk-Nusair) parameterization
is standard in most codes, giving $\varepsilon_{\mathrm{xc}}$ as a
function of $r_s$ with both exchange ($\propto r_s^{-1}$, Dirac 1930)
and correlation (QMC) contributions.
Because each point $\bm{r}$ contributes independently (no nonlocal coupling),
the functional derivative reduces to an ordinary derivative of the integrand
with respect to $\rho$:

\begin{derivbox}
Perturb $\rho\to\rho+\varepsilon\phi$:
\begin{align}
  \Exc^{\mathrm{LDA}}[\rho+\varepsilon\phi]
  &= \int \varepsilon_{\mathrm{xc}}\!\left(\rho(\bm{r})+\varepsilon\phi(\bm{r})\right)
    \bigl(\rho(\bm{r})+\varepsilon\phi(\bm{r})\bigr)\,\dd^3r \notag\\
  &= \int
    \Bigl[\varepsilon_{\mathrm{xc}}(\rho)
    + \varepsilon\phi\,\varepsilon_{\mathrm{xc}}'(\rho)
    + O(\varepsilon^2)\Bigr]
    \bigl[\rho + \varepsilon\phi\bigr]\,\dd^3r \notag\\
  &= \Exc^{\mathrm{LDA}}[\rho]
    + \varepsilon\int\phi(\bm{r})
      \underbrace{\left[\varepsilon_{\mathrm{xc}}(\rho)
        + \rho\,\varepsilon_{\mathrm{xc}}'(\rho)\right]}_{
        = \dd[\rho\,\varepsilon_{\mathrm{xc}}]/\dd\rho}
      \,\dd^3r + O(\varepsilon^2).
\end{align}
Reading off the coefficient of $\phi(\bm{r})$:
\begin{equation}
  \boxed{\Vxc^{\mathrm{LDA}}(\bm{r})
    = \frac{\delta\Exc^{\mathrm{LDA}}}{\delta\rho(\bm{r})}
    = \frac{\dd\bigl[\rho\,\varepsilon_{\mathrm{xc}}(\rho)\bigr]}{\dd\rho}
    \bigg|_{\rho = \rho(\bm{r})}.}
  \label{eq:vxc_lda}
\end{equation}
\end{derivbox}

\noindent\textbf{Key message.}
For LDA, the functional derivative is an ordinary derivative of the integrand
because LDA is \emph{local in space}: the XC energy density at $\bm{r}$
depends only on $\rho(\bm{r})$, not on $\rho(\bm{r}')$ at any other point.
For GGAs the same procedure applies but one must integrate by parts
because $\Exc^{\mathrm{GGA}}$ depends on $|\nabla\rho|$ as well.
Here is the essential step.
A GGA energy density has the form
$\varepsilon_{\mathrm{xc}}(\rho,|\nabla\rho|)$, so:
\begin{equation}
  \Exc^{\mathrm{GGA}}[\rho]
  = \int \varepsilon_{\mathrm{xc}}\!\bigl(\rho(\bm{r}),|\nabla\rho(\bm{r})|\bigr)
    \,\rho(\bm{r})\,\dd^3r.
\end{equation}
Perturbing $\rho\to\rho+\varepsilon\phi$ generates a term
$\varepsilon_{\mathrm{xc},s}\,(\nabla\rho\cdot\nabla\phi)/|\nabla\rho|$
(where $\varepsilon_{\mathrm{xc},s} = \partial(\rho\varepsilon_{\mathrm{xc}})/\partial s$,
$s=|\nabla\rho|$) that contains $\nabla\phi$ rather than $\phi$.
To read off the coefficient of $\phi$ one integrates by parts,
using $\int f\,\nabla\phi\,\dd^3r = -\int(\nabla f)\phi\,\dd^3r$
(boundary term vanishes for localised densities):
\begin{equation}
  \Vxc^{\mathrm{GGA}}(\bm{r})
  = \frac{\partial(\rho\,\varepsilon_{\mathrm{xc}})}{\partial\rho}
  - \nabla\cdot\left[
      \frac{\partial(\rho\,\varepsilon_{\mathrm{xc}})}{\partial|\nabla\rho|}
      \frac{\nabla\rho}{|\nabla\rho|}
    \right].
  \label{eq:vxc_gga}
\end{equation}
The second term is the gradient correction absent from LDA.
In practice, codes evaluate both terms on the quadrature grid via
automatic differentiation or analytical GGA implementations.

\begin{insight}
\textbf{How $V_\text{xc}^\text{LDA}$ enters the KS matrix in a Gaussian-basis code.}
$V_\text{xc}^\text{LDA}(\bm{r}) = \dd[\rho\varepsilon_{\rm xc}]/\dd\rho|_{\rho(\bm{r})}$
is a \emph{multiplicative} (local) potential: it is just a number at each
point $\bm{r}$, not an integral operator.
In a Gaussian basis $\{\chi_\mu\}$, it enters the KS Fock matrix via:
\begin{equation}
  (V_{\rm xc})_{\mu\nu}
  = \int \chi_\mu(\bm{r})\,V_\text{xc}^{\rm LDA}(\bm{r})\,\chi_\nu(\bm{r})
    \,\dd^3r
  \approx \sum_g w_g\,\chi_\mu(\bm{r}_g)\,V_\text{xc}^{\rm LDA}(\rho(\bm{r}_g))
    \,\chi_\nu(\bm{r}_g),
  \label{eq:vxc_matrix}
\end{equation}
where $\{(\bm{r}_g, w_g)\}$ is a numerical quadrature grid (e.g.\ Becke
grid in molecular codes, FFT grid in plane-wave codes).
This is the computational realization of Eq.~(\ref{eq:vxc_lda}): the
functional derivative is evaluated at each grid point as an ordinary
number, then integrated against basis-function pairs.

In RMG (real-space grid), each grid point \emph{is} a quadrature point,
so the sum over $g$ is identical to the matrix element:
$(V_{\rm xc})_{ij} = \sum_g w_g\,\phi_i(\bm{r}_g)\,V_\text{xc}(\bm{r}_g)\,\phi_j(\bm{r}_g)$.
The two codes differ only in the basis set, not in how $V_\text{xc}$ is applied.
\end{insight}

% ----------------------------------------------------------------------------
\subsection{The Second Functional Derivative and the XC Kernel}
% ----------------------------------------------------------------------------

Just as ordinary second derivatives $f''(x) = \dd^2f/\dd x^2$ measure the
\emph{curvature} of $f$, the second functional derivative
$\delta^2 F/\delta f(x)\,\delta f(x')$ measures the \emph{nonlocal curvature}
of a functional.

\paragraph{Definition.}
The second functional derivative is defined iteratively: differentiate
$\delta F/\delta f(x)$ --- which is itself a functional of $f$ --- with
respect to $f$ at a second point $x'$:
\begin{equation}
  \frac{\delta^2 F}{\delta f(x)\,\delta f(x')}
  \;=\;
  \frac{\delta}{\delta f(x')}
  \left[\frac{\delta F}{\delta f(x)}\right].
  \label{eq:sfd_def}
\end{equation}

\paragraph{Gradient analogy.}
On a finite grid, the first derivative is the gradient vector
$g_i = \partial F/\partial a_i$.
The second functional derivative is the \emph{Hessian matrix}:
\begin{equation}
  \frac{\delta^2 F}{\delta f(x)\,\delta f(x')}
  \;\longleftrightarrow\;
  H_{ij} = \frac{\partial^2 F}{\partial a_i\,\partial a_j}.
  \label{eq:sfd_hessian}
\end{equation}
The Hessian $H_{ij}$ tells you how the gradient component at grid point $i$
changes when you change the value at grid point $j$.
The second functional derivative $\delta^2 F/\delta f(x)\,\delta f(x')$ is
the continuum limit: how the ``gradient'' $\delta F/\delta f(x)$ changes when
$f$ is nudged at a different point $x'$.

\paragraph{Physical meaning for DFT.}
For the XC energy, the second functional derivative is the \emph{XC kernel}:
\begin{equation}
  \fxc^{\mathrm{gs}}(\bm{r},\bm{r}')
  = \frac{\delta^2\Exc}{\delta\rho(\bm{r})\,\delta\rho(\bm{r}')}
  = \frac{\delta\Vxc(\bm{r})}{\delta\rho(\bm{r}')}.
  \label{eq:fxc_kernel_primer}
\end{equation}

\noindent\textbf{Reading this equation:}
$\fxc(\bm{r},\bm{r}')$ tells you how much the XC potential at point $\bm{r}$
changes when the electron density is nudged by a tiny amount at the
\emph{different} point $\bm{r}'$.
If $\bm{r}\neq\bm{r}'$, this measures the \emph{spatial nonlocality} of the XC
interaction: an electron added at $\bm{r}'$ ``feels'' the XC environment change
at $\bm{r}$.

\paragraph{Worked example: second derivative of $F[f]=\int f^2\,\dd x$.}
We already know from Worked Example~1 that $\delta F/\delta f(x) = 2f(x)$.
Apply the definition~(\ref{eq:sfd_def}): differentiate $2f(x)$ --- which is
a local functional of $f$ --- with respect to $f$ at the second point $x'$.

\begin{derivbox}
By the composite chain rule~(\ref{eq:chain_rule_local}) with $g(f) = 2f$
and $g'(f) = 2$:
\begin{equation}
  \frac{\delta^2 F}{\delta f(x)\,\delta f(x')}
  = \frac{\delta}{\delta f(x')}\bigl[2f(x)\bigr]
  = 2\,\delta(x - x').
  \label{eq:sfd_example}
\end{equation}
\textbf{Grid check.}
On a grid with spacing $h$, $F \approx \sum_i f_i^2 h$, so
$\partial F/\partial f_i = 2f_i h$ and
$\partial^2 F/\partial f_i\,\partial f_j = 2h\,\delta_{ij}$.
Dividing by $h^2$ to convert two grid-value derivatives to continuum:
$2h\,\delta_{ij}/h^2 = 2\delta_{ij}/h \to 2\delta(x-x')$ as $h\to 0$.
Consistent.
\end{derivbox}

\noindent\textbf{Physical meaning of the result.}
The Hessian of $F[f]=\int f^2\,\dd x$ is \emph{diagonal} in the continuum
sense: nudging $f$ at $x'$ only changes the gradient $\delta F/\delta f(x)$
at the \emph{same} point $x = x'$, by the amount $2$.
A nonlocal functional would have off-diagonal structure: nudging at $x'$
would affect the gradient at distant $x\neq x'$.
For DFT: LDA is diagonal ($\fxc \propto \delta(\bm{r}-\bm{r}')$);
hybrid functionals and exact exchange are off-diagonal.

\begin{keyresult}[title={Why $\fxc$ Controls Linear Response}]
Suppose the system is in its ground state and is then perturbed by
$\delta V_\text{ext}(\bm{r},\omega)$.
The induced density change is:
\begin{equation}
  \delta\rho(\bm{r},\omega)
  = \int \chi(\bm{r},\bm{r}',\omega)\,\delta V_\text{ext}(\bm{r}',\omega)\,\dd^3r',
  \label{eq:chi_def}
\end{equation}
where $\chi$ is the interacting density-response function.

The key idea is that the KS electrons respond to the \emph{total} effective
potential $V_\text{eff} = V_\text{ext} + V_H + V_\text{xc}$, not just
$V_\text{ext}$ alone.  Because $V_H$ and $V_\text{xc}$ themselves depend on
$\delta\rho$, the perturbation is self-consistently amplified.
Writing $\delta V_\text{eff} = \delta V_\text{ext} + (1/|\bm{r}-\bm{r}'| + \fxc)\delta\rho$
and using $\delta\rho = \chi_s\,\delta V_\text{eff}$ (KS electrons respond
to $V_\text{eff}$) gives a self-consistent equation for $\delta\rho$.
Substituting $\chi\,\delta V_\text{ext} \equiv \delta\rho$ to eliminate
$\delta V_\text{ext}$ yields the Dyson equation for the interacting response
$\chi$ in terms of the non-interacting KS response $\chi_s$:
\begin{equation}
  \chi(\bm{r},\bm{r}',\omega) = \chi_s(\bm{r},\bm{r}',\omega)
  + \iint \chi_s(\bm{r},\bm{r}'',\omega)
    \left[\frac{1}{|\bm{r}''-\bm{r}'''|}
    + \fxc(\bm{r}'',\bm{r}''',\omega)\right]
    \chi(\bm{r}''',\bm{r}',\omega)\,\dd^3r''\,\dd^3r'''.
  \label{eq:dyson}
\end{equation}
The XC kernel $\fxc$ enters as the \emph{exchange-correlation contribution
to the effective screening}.
In the adiabatic approximation, $\fxc(\omega) = \fxc^{\mathrm{gs}}$
is frequency-independent.
\end{keyresult}

\begin{keyresult}[title={Composite Chain Rule for Local Functionals}]
Let $g(\rho)$ be an ordinary differentiable function of a real variable, and
define the functional $G[\rho]$ by pointwise evaluation:
$G[\rho](\bm{r}) \equiv g\!\left(\rho(\bm{r})\right)$.
Then:
\begin{equation}
  \frac{\delta G[\rho](\bm{r})}{\delta\rho(\bm{r}')}
  = g'\!\left(\rho(\bm{r})\right)\,\delta(\bm{r}-\bm{r}').
  \label{eq:chain_rule_local}
\end{equation}
\begin{derivbox}[title={Proof}]
Perturb $\rho \to \rho + \varepsilon\,\delta(\cdot - \bm{r}')$ (a spike at $\bm{r}'$):
\begin{equation}
  G[\rho + \varepsilon\,\delta(\cdot-\bm{r}')\,](\bm{r})
  = g\!\left(\rho(\bm{r}) + \varepsilon\,\delta(\bm{r}-\bm{r}')\right).
\end{equation}
Taylor-expand $g$ to first order in $\varepsilon$:
\begin{equation}
  = g\!\left(\rho(\bm{r})\right)
    + \varepsilon\,g'\!\left(\rho(\bm{r})\right)\,\delta(\bm{r}-\bm{r}')
    + O(\varepsilon^2).
\end{equation}
Subtract $G[\rho](\bm{r})$, divide by $\varepsilon$, take $\varepsilon\to 0$:
\begin{equation}
  \frac{\delta G[\rho](\bm{r})}{\delta\rho(\bm{r}')}
  = g'\!\left(\rho(\bm{r})\right)\,\delta(\bm{r}-\bm{r}'). \qquad\square
\end{equation}
\end{derivbox}
\noindent\textbf{Interpretation.}
Because $g$ acts locally (it only sees $\rho(\bm{r})$, not $\rho$ at any other
point), nudging the density at $\bm{r}'$ can only affect $G(\bm{r})$ when
$\bm{r}' = \bm{r}$.  The Dirac delta enforces this: the sensitivity is
$g'(\rho(\bm{r}))$ at $\bm{r}' = \bm{r}$ and zero everywhere else.
This is precisely why LDA produces a spatially local kernel.
\end{keyresult}

\paragraph{Computing $\fxc$ for LDA.}
Since $\Vxc^{\mathrm{LDA}}(\bm{r}) = v_\text{xc}(\rho(\bm{r}))$ is a local function
of $\rho$, Eq.~(\ref{eq:chain_rule_local}) applies with $g = v_\text{xc}$:
\begin{derivbox}
\begin{align}
  \fxc^{\mathrm{LDA}}(\bm{r},\bm{r}')
  &= \frac{\delta\Vxc^{\mathrm{LDA}}(\bm{r})}{\delta\rho(\bm{r}')} \notag\\
  &= \frac{\delta}{\delta\rho(\bm{r}')}\left[
    \frac{\dd[\rho\,\varepsilon_{\mathrm{xc}}(\rho)]}{\dd\rho}
    \bigg|_{\rho(\bm{r})}
    \right] \notag\\
  &\overset{(\ref{eq:chain_rule_local})}{=}
    \frac{\dd^2[\rho\,\varepsilon_{\mathrm{xc}}(\rho)]}{\dd\rho^2}
    \bigg|_{\rho(\bm{r})}
    \cdot\delta(\bm{r}-\bm{r}').
  \label{eq:fxc_lda_full}
\end{align}
The middle line identifies $g = \dd[\rho\varepsilon_{\mathrm{xc}}]/\dd\rho$,
an ordinary function of $\rho$; the chain rule then supplies the factor
$g' = \dd^2[\rho\varepsilon_{\mathrm{xc}}]/\dd\rho^2$ multiplied by
$\delta(\bm{r}-\bm{r}')$.
\end{derivbox}

\noindent The delta function $\delta(\bm{r}-\bm{r}')$ encodes the fact that
LDA is \emph{spatially local}: nudging the density at $\bm{r}'$ only changes
$\Vxc$ at that same point $\bm{r}' = \bm{r}$.
A nonlocal functional (e.g.\ a hybrid with exact exchange) would produce a
kernel that is nonzero for $\bm{r}\neq\bm{r}'$.

% ----------------------------------------------------------------------------
\subsection{Summary Table: Analogy Between Ordinary and Functional Calculus}
% ----------------------------------------------------------------------------

\begin{center}
\renewcommand{\arraystretch}{1.5}
\begin{tabular}{lll}
\toprule
\textbf{Ordinary calculus} & & \textbf{Functional calculus} \\
\midrule
Function $f(x):\mathbb{R}\to\mathbb{R}$ & $\longleftrightarrow$ &
  Functional $F[f]:\{\text{functions}\}\to\mathbb{R}$ \\
Derivative $f'(x_0)$ & $\longleftrightarrow$ &
  Functional deriv.\ $\delta F/\delta f(x_0)$ \\
Gradient $\partial F/\partial a_i$ & $\longleftrightarrow$ &
  $\delta F/\delta f(x)$ at continuum point $x$ \\
Hessian $\partial^2 F/\partial a_i\partial a_j$ & $\longleftrightarrow$ &
  $\delta^2 F/\delta f(x)\,\delta f(x')$ \\
Chain rule: $\dd F/\dd x = \sum_i (\partial F/\partial a_i)(\dd a_i/\dd x)$ &
  $\longleftrightarrow$ &
  $\delta F/\delta g(y) = \int (\delta F/\delta f(x))(\delta f(x)/\delta g(y))\,\dd x$ \\
Local function: $F(a_i)$ depends only on $a_i$ & $\longleftrightarrow$ &
  Local functional: $F[f]$ at $x$ depends only on $f(x)$ (e.g.\ LDA) \\
$\partial^2 F/\partial a_i\partial a_j = 0$ for $i\neq j$ (diagonal Hessian) &
  $\longleftrightarrow$ &
  $\fxc^{\mathrm{LDA}}(\bm{r},\bm{r}') \propto \delta(\bm{r}-\bm{r}')$ \\
\bottomrule
\end{tabular}
\renewcommand{\arraystretch}{1}
\end{center}

\paragraph{The functional chain rule: when does it arise?}
The chain rule row in the table applies whenever one functional depends on
another.
$f$ is a functional of $g$ when $f[g]$ maps functions $g$ to functions $f$:
each value $f(x)$ is determined by the entire function $g(\cdot)$, not just
$g$ at a single point.

\textbf{When this arises in DFT.}
The density $\rho$ is a functional of the external potential $V_\text{ext}$
(via the ground-state KS equations), so any functional of $\rho$ is
implicitly a functional of $V_\text{ext}$.
The chain rule then relates $\delta F/\delta V_\text{ext}$ to
$\delta F/\delta\rho$ and $\delta\rho/\delta V_\text{ext} \equiv \chi_s$
(the KS response function).

\textbf{Concrete example.}
Let $f(x) = \int K(x,y)\,g(y)\,\dd y$ (a linear transform of $g$) and
$F[f] = \int f(x)^2\,\dd x$.
We know $\delta F/\delta f(x) = 2f(x)$ and
$\delta f(x)/\delta g(y) = K(x,y)$.
The chain rule gives:
\begin{equation}
  \frac{\delta F}{\delta g(y)}
  = \int \frac{\delta F}{\delta f(x)}\,\frac{\delta f(x)}{\delta g(y)}\,\dd x
  = \int 2f(x)\,K(x,y)\,\dd x
  = 2\int K(x,y)\,f(x)\,\dd x.
  \label{eq:chain_rule_example}
\end{equation}
\textbf{Check by direct calculation.}
$F[g] = \int\!\bigl(\int K(x,y)g(y)\,\dd y\bigr)^2\!\dd x$.
Perturb $g\to g+\varepsilon\phi$ and take $\varepsilon\to 0$:
$\delta F/\delta g(y) = 2\int K(x,y)f(x)\,\dd x$. Consistent. $\checkmark$

\begin{insight}
\textbf{Why $\fxc^{\mathrm{LDA}}$ is a delta function in space but not in time
(for TDDFT).}
Spatially, LDA is local: $\Vxc(\bm{r})$ depends only on $\rho(\bm{r})$,
not on $\rho(\bm{r}')$ at other points.
This makes $\delta\Vxc(\bm{r})/\delta\rho(\bm{r}') \propto \delta(\bm{r}-\bm{r}')$.
The adiabatic approximation further makes it a delta function \emph{in time}:
$\fxc(\bm{r},\bm{r}',t-t') = \fxc^{\mathrm{gs}}(\bm{r},\bm{r}')\,\delta(t-t')$.
Both delta functions together say:
the XC response has zero range in space and zero memory in time.
This is what Section~\ref{sec:memory} now analyzes in detail.
\end{insight}

% ============================================================================
\section{Memory: What Is Discarded and Why Derivatives Are Insufficient}
\label{sec:memory}
% ============================================================================

\subsection{Memory as a Convolution Integral}

\begin{keyresult}[title={Functional Taylor Expansion}]
Any sufficiently smooth functional $F[f]$ can be expanded around a
reference $f_0$ in powers of the deviation $\delta f = f - f_0$:
\begin{equation}
  F[f_0 + \delta f]
  = F[f_0]
    + \int \frac{\delta F}{\delta f(x)}\bigg|_{f_0}\!\!\delta f(x)\,\dd x
    + \frac{1}{2}\iint
        \frac{\delta^2 F}{\delta f(x)\,\delta f(x')}\bigg|_{f_0}\!\!
        \delta f(x)\,\delta f(x')\,\dd x\,\dd x'
    + \cdots
  \label{eq:func_taylor}
\end{equation}
This is the functional analogue of the ordinary Taylor series
$F(x_0+\delta x) = F(x_0) + F'(x_0)\delta x + \tfrac{1}{2}F''(x_0)(\delta x)^2 + \cdots$,
with the sum over discrete terms replaced by integrals over continuous variables.
Truncating at first order gives the \textbf{linear response} approximation:
\begin{equation}
  \delta F \equiv F[f_0+\delta f] - F[f_0]
  \approx \int \frac{\delta F}{\delta f(x)}\bigg|_{f_0}\!\!\delta f(x)\,\dd x.
  \label{eq:linear_response}
\end{equation}
\end{keyresult}

To see concretely what ``memory'' means, linearize $\Vxc$ around a
reference density $\rho_0$ using Eq.~(\ref{eq:linear_response}). For a small perturbation
$\delta\rho(\bm{r},t)$:
\begin{equation}
  \delta\Vxc(\bm{r},t)
  = \int_0^t\!\dd t'\!\int\!\dd^3r'\;
    \underbrace{\frac{\delta\Vxc(\bm{r},t)}{\delta\rho(\bm{r}',t')}}_{\displaystyle\fxc(\bm{r},\bm{r}',t-t')}
    \delta\rho(\bm{r}',t').
  \label{eq:fxc_def}
\end{equation}
The quantity $\fxc(\bm{r},\bm{r}',\tau)$ is the \emph{XC kernel}.
\textbf{It is exactly the second functional derivative of the XC energy
computed in Section~\ref{sec:funcderiv}}, now promoted to a time-dependent
object:
\begin{equation}
  \fxc(\bm{r},\bm{r}',t-t')
  = \frac{\delta\Vxc(\bm{r},t)}{\delta\rho(\bm{r}',t')}
  = \frac{\delta^2\Exc[\rho]}{\delta\rho(\bm{r},t)\,\delta\rho(\bm{r}',t')}.
  \label{eq:fxc_is_sfd}
\end{equation}
In the ground-state (static) limit $t\to t'$, this reduces to
$\fxc^{\rm gs}(\bm{r},\bm{r}') = \delta^2\Exc/\delta\rho(\bm{r})\delta\rho(\bm{r}')$
from Eq.~(\ref{eq:fxc_kernel_primer}).
The memory structure ($\tau = t-t'$ dependence) is the new feature:
it encodes how a density perturbation at time $t'$ continues to affect
the XC potential at all later times $t > t'$, i.e., XC has a causal
memory kernel. $\fxc = 0$ for $\tau < 0$ (causality).

\begin{insight}
The exact $\Vxc$ is a \textbf{convolution over history}: the XC
response today depends on the whole trajectory $\delta\rho(\bm{r}',t')$
for all past times $t'\leq t$, weighted by the memory kernel
$\fxc(\tau)$.  This is analogous to a viscoelastic material whose
stress today reflects the full history of its deformation --- not just
its current state.
\end{insight}

The adiabatic approximation drops the time-nonlocality entirely.
In the adiabatic approximation, $V_\text{xc}(\bm{r},t) = V_\text{xc}^{\rm gs}[\rho(\cdot,t)](\bm{r})$
depends only on the density \emph{at the same instant} $t$, not on any
earlier time $t'$.
Therefore, varying with respect to $\rho(\bm{r}',t')$:
\begin{itemize}[noitemsep]
  \item For $t' = t$: the variation gives the \emph{static} XC kernel
    $\fxc^{\rm gs}(\bm{r},\bm{r}') = \delta V_\text{xc}^{\rm gs}/\delta\rho(\bm{r}')$,
    exactly the ground-state second functional derivative.
  \item For $t' \neq t$: $V_\text{xc}(t)$ does not depend on $\rho(t')$
    at all, so $\delta V_\text{xc}(\bm{r},t)/\delta\rho(\bm{r}',t') = 0$.
\end{itemize}
Combined, the adiabatic kernel is:
\begin{equation}
  \fxc^{\mathrm{ATDDFT}}(\bm{r},\bm{r}',\tau)
  = \fxc^{\mathrm{gs}}(\bm{r},\bm{r}')\;\delta(\tau),
  \label{eq:fxc_adiab}
\end{equation}
a delta function in time: the response depends only on the density
at the same instant, with zero memory.

\subsection{The Frequency-Domain Picture}

Causality — $\fxc(\tau) = 0$ for $\tau < 0$ — means the standard
two-sided Fourier transform $\int_{-\infty}^{\infty}\fxc(\tau)e^{i\omega\tau}\,\dd\tau$
reduces to a one-sided integral: the $\tau<0$ half contributes nothing.
This one-sided transform is the \textbf{Laplace transform evaluated on
the real axis} (or equivalently, the Fourier transform restricted to
causal functions):
\begin{equation}
  \fxc(\bm{r},\bm{r}',\omega)
  = \int_0^\infty\!\fxc(\bm{r},\bm{r}',\tau)\,e^{\ii\omega\tau}\,\dd\tau.
  \label{eq:fxc_fourier}
\end{equation}
The sign convention $e^{+i\omega\tau}$ (rather than $e^{-i\omega\tau}$)
is chosen so that $\fxc(\omega)$ is analytic in the upper complex
half-plane $\operatorname{Im}\omega > 0$, which is the standard
physics convention for retarded response functions and is required for
the Kramers-Kronig relations (Section~\ref{sec:cdft}).
The adiabatic kernel corresponds to:
\begin{equation}
  \fxc^{\mathrm{ATDDFT}}(\omega) = \fxc^{\mathrm{gs}} = \text{const.}
  \quad\text{(independent of }\omega\text{)}.
  \label{eq:fxc_static}
\end{equation}
The \emph{exact} $\fxc(\omega)$ has a much richer structure.

\paragraph{Why response functions have poles: the resonance picture.}
Consider driving a quantum system with a perturbation at frequency $\omega$.
The response $\delta\rho(\omega)$ is proportional to the resolvent
$(H - \omega)^{-1}$, which diverges whenever $\omega$ equals an eigenvalue of $H$.
In Dirac notation, the exact density-response function has the
\textbf{Lehmann representation}:
\begin{equation}
  \chi(\bm{r},\bm{r}',\omega)
  = \sum_k \left[
      \frac{\langle 0|\hat\rho(\bm{r})|k\rangle\langle k|\hat\rho(\bm{r}')|0\rangle}
           {\omega - \Omega_k + \ii\eta}
    - \frac{\langle 0|\hat\rho(\bm{r}')|k\rangle\langle k|\hat\rho(\bm{r})|0\rangle}
           {\omega + \Omega_k + \ii\eta}
    \right],
  \label{eq:lehmann}
\end{equation}
where $\Omega_k = E_k - E_0 > 0$ are the exact excitation energies and
$|k\rangle$ the exact excited states ($\eta\to 0^+$ enforces causality).
Each term $1/(\omega - \Omega_k)$ is a \emph{simple pole} at $\omega = \Omega_k$:
a sharp divergence at the resonance frequency.
Physically, near $\omega\approx\Omega_k$ the perturbation drives the system
at exactly the frequency where it ``wants to oscillate,'' producing an
unbounded linear response (bounded in practice by damping or finite pulse duration).

Because $\fxc$ is defined via the Dyson equation~(\ref{eq:dyson})
relating $\chi_s$ (KS, poles at KS excitation energies) to $\chi$
(interacting, poles at exact excitation energies $\Omega_k$), $\fxc$
must itself have poles at $\Omega_k$ to shift the KS poles to the correct
positions.  For a finite system:
\begin{equation}
  \fxc(\omega) \sim \sum_k \frac{R_k}{\omega - \Omega_k}
  + \text{regular terms},
  \label{eq:fxc_poles}
\end{equation}
where $R_k$ are residues encoding the coupling strengths.  The imaginary
part $\operatorname{Im}\fxc(\omega)\neq 0$ for $\omega>0$ encodes dissipation.

\subsection{Why Time Derivatives Cannot Replace Full Memory}
\label{sec:derivatives}

A natural attempt to go beyond the adiabatic limit is to include
time derivatives of the density:
\begin{equation}
  \Vxc[\rho](t) \approx \Vxc^{\mathrm{gs}}[\rho(t)]
  + \alpha_1\frac{\partial\rho}{\partial t}
  + \alpha_2\frac{\partial^2\rho}{\partial t^2} + \cdots
  \label{eq:gradient_expansion}
\end{equation}
This is the \emph{adiabatic gradient expansion in time}.

In frequency space, including terms up to order $N$ corresponds to
replacing the full $\fxc(\omega)$ by a polynomial of degree $N$:
\begin{equation}
  \fxc^{\mathrm{ATDDFT}+N}(\omega)
  \approx c_0 + c_1\omega + c_2\omega^2 + \cdots + c_N\omega^N.
  \label{eq:polynomial}
\end{equation}

\begin{keyresult}
A polynomial in $\omega$ cannot reproduce a function with poles.
No finite-order derivative expansion can capture the poles in
Eq.~(\ref{eq:fxc_poles}).  The poles are essential for:
\begin{enumerate}[noitemsep]
  \item Exact resonance positions of excitation energies.
  \item Finite linewidths (imaginary part of poles).
  \item Double excitations (poles at two-electron excitation frequencies).
\end{enumerate}
\end{keyresult}

\begin{derivbox}
\textbf{Explicit demonstration.}
Suppose the exact $\fxc(\omega)$ has a single pole at $\Omega_1$:
\[
  \fxc(\omega) = \frac{R_1}{\omega - \Omega_1}.
\]
Its Taylor expansion around $\omega=0$ is:
\[
  \fxc(\omega) = -\frac{R_1}{\Omega_1}\sum_{n=0}^\infty
  \left(\frac{\omega}{\Omega_1}\right)^n
  = -\frac{R_1}{\Omega_1} - \frac{R_1}{\Omega_1^2}\omega
  - \frac{R_1}{\Omega_1^3}\omega^2 - \cdots
\]
This series \emph{diverges} for $|\omega|>|\Omega_1|$ and never
converges to the correct form at $\omega=\Omega_1$.  The pole cannot
be represented by any finite polynomial.

\textbf{Deeper reason (complex analysis).}
A polynomial $p(\omega) = \sum_{n=0}^N c_n\omega^n$ is an \emph{entire
function}: it is analytic everywhere in the complex $\omega$-plane and
has no singularities at any finite point.
A function with a pole at $\Omega_1$ is \emph{not} entire --- it is
meromorphic, with a singularity at $\Omega_1$.
No entire function can equal a meromorphic function:
their analyticity structures are fundamentally different.
Increasing $N$ does not help; adding more polynomial terms keeps the
function entire.  The mismatch is not a convergence-radius problem but
a topological one: you cannot continuously deform an analytic function
into one with a pole.

Physically: near $\omega \approx \Omega_1$, the electron system has a
resonant response that persists in time.  The polynomial approximation
has no singularity and therefore no resonant buildup.
\end{derivbox}

\begin{remark}
Driving with $\cos(\omega_{\rm drive}\,t)$ at the resonance frequency
does produce a resonant buildup in the \emph{total} density response
even in adiabatic TDDFT, because the KS system has its own resonance
near $\omega_0$.  What the adiabatic approximation gets wrong are:
(a) the exact position of $\omega_0$ (determined by the poles of
$\fxc$); (b) the linewidth (Im~$\fxc = 0$ in the adiabatic limit);
(c) states dominated by double-excitation character (Section~\ref{sec:doubles}).
\end{remark}

% ============================================================================
\section{Dissipation, Unitarity, and Dephasing}
\label{sec:dissipation}
% ============================================================================

\subsection{The Apparent Paradox}

Adiabatic TDDFT uses a Hermitian KS Hamiltonian $\Hks^\dagger = \Hks$.
The KS one-body density matrix is $P_{ij}(t) = \langle\psi_j|\psi_i\rangle$
(more precisely, for a set of occupied orbitals,
$P = \sum_n f_n |\psi_n\rangle\langle\psi_n|$).
Differentiating with respect to time and using the TDKS equation
$\ii\partial_t|\psi_n\rangle = \Hks|\psi_n\rangle$ (and its conjugate
$-\ii\partial_t\langle\psi_n| = \langle\psi_n|\Hks$):
\begin{align}
  \ii\dot{P}
  &= \ii\sum_n f_n\bigl(\partial_t|\psi_n\rangle\langle\psi_n|
                       + |\psi_n\rangle\partial_t\langle\psi_n|\bigr)
   = \sum_n f_n\bigl(\Hks|\psi_n\rangle\langle\psi_n|
                     - |\psi_n\rangle\langle\psi_n|\Hks\bigr)
   = \comm{\Hks}{P}.
  \label{eq:vonneumann_derive}
\end{align}
This is the \textbf{von Neumann equation} $\ii\dot{P} = \comm{\Hks}{P}$.

\begin{insight}
\textbf{MO-basis vs AO-basis density matrix: connecting RMG to Gaussian codes.}
RMG propagates $P(t)$ in the basis of \emph{ground-state KS orbitals}
$\{\phi_n(\bm{r})\}$ (MO basis), so $P_{ij}(t) = \langle\phi_j|P(t)|\phi_i\rangle$
is a matrix in MO space.

Gaussian codes such as NWChem or Gaussian propagate the density matrix in
the \emph{atomic orbital} (AO) basis $\{\chi_\mu(\bm{r})\}$:
\begin{equation}
  P_{\mu\nu}(t) = \langle\chi_\nu|P(t)|\chi_\mu\rangle
  = \sum_{kl} C_{\mu k}^*\,P_{kl}(t)\,C_{\nu l},
  \label{eq:P_AO_MO}
\end{equation}
where $C_{\mu n}$ is the MO coefficient matrix ($\phi_n = \sum_\mu C_{\mu n}\chi_\mu$).
In matrix notation: $\bm{P}_\text{AO} = \bm{C}\bm{P}_\text{MO}\bm{C}^\dagger$.

The von Neumann equation in the AO basis includes the overlap matrix
$S_{\mu\nu} = \langle\chi_\mu|\chi_\nu\rangle$ (non-orthogonal basis):
\begin{equation}
  \ii\bm{S}\dot{\bm{P}}_\text{AO} = \bm{H}_\text{AO}\bm{P}_\text{AO}\bm{S}
                                   - \bm{S}\bm{P}_\text{AO}\bm{H}_\text{AO}.
\end{equation}
In the MO basis $\bm{S} = \bm{1}$ (orthonormal KS orbitals), so the
equation reduces to the clean form $\ii\dot{P} = [H,P]$ used in RMG.
The two representations are physically identical; the MO basis simply
removes the overhead of carrying $\bm{S}$.
\end{insight}

The \textbf{solution} of the von Neumann equation over one timestep $\Delta t$
is given by the \textbf{Magnus propagator}.
The equation of motion $\ii\dot{P} = [H,P]$ is an operator differential
equation; its formal solution for a time-dependent $H$ is \emph{not} simply
$P(t) = e^{-iHt}P(0)e^{+iHt}$ (that would only be correct for constant $H$).
Instead, the Magnus expansion provides the exact solution as
$P(t+\Delta t) = e^{-i\Omega}P(t)e^{+i\Omega}$, where $\Omega$ is
a Hermitian operator built from time-ordered integrals of $H$.
To second order (the level used in RMG):
\begin{equation}
  P(t+\Delta t) = e^{-\ii\Omega}\,P(t)\,e^{+\ii\Omega}, \qquad
  \Omega \approx \tfrac{1}{2}(H(t)+H(t+\Delta t))\,\Delta t,
  \quad \Omega = \Omega^\dagger.
  \label{eq:magnus}
\end{equation}
This is a unitary similarity transform: $e^{-i\Omega}$ is the propagator,
and its Hermitian conjugate $e^{+i\Omega}$ appears because $P$ carries
a ket on the left and a bra on the right ($P = |\psi\rangle\langle\psi|$).  One might then ask: can the
exact TDDFT functional (with memory) restore dissipation?  And does
unitarity preclude all relaxation?

The answer requires distinguishing three different things:
the $N$-body wavefunction, the one-body density matrix, and what
it means to have a ``closed'' system.

\subsection{Unitarity of the Magnus Propagator and Conservation of Occupations}

\begin{derivbox}
\textbf{Eigenvalues of $P$ are conserved under unitary propagation.}

Let $P(t) = V\Lambda V^\dagger$ be the eigendecomposition of $P$ at time $t$,
with eigenvalues $\Lambda = \mathrm{diag}(f_1, f_2, \ldots)$ (occupation numbers).
Under the Magnus step Eq.~(\ref{eq:magnus}):
\begin{align}
  P(t+\Delta t)
  &= e^{-\ii\Omega}V\Lambda V^\dagger e^{+\ii\Omega} \notag \\
  &= \underbrace{\left(e^{-\ii\Omega}V\right)}_{\tilde{V}}
     \;\Lambda\;
     \underbrace{\left(e^{-\ii\Omega}V\right)^\dagger}_{\tilde{V}^\dagger}.
\end{align}
Since $e^{-\ii\Omega}$ is unitary (because $\Omega$ is Hermitian),
$\tilde{V}$ is unitary, and $P(t+\Delta t) = \tilde{V}\Lambda\tilde{V}^\dagger$
has the \emph{same eigenvalues} $\Lambda$.

Consequence: the occupation numbers $\{f_n\}$ are constants of motion in
KS-TDDFT.  No population can transfer from an excited state to the
ground state --- ever.  This is exact regardless of whether
$V_{\rm xc}$ has memory.
\end{derivbox}

Therefore: \textbf{T$_1$ relaxation (population decay) is absent in
TDDFT}, not because of the adiabatic approximation, but because the
KS Hamiltonian is Hermitian and the evolution is unitary.

\subsection{T\textsubscript{1} and T\textsubscript{2}: The Bloch Equation Analogy}

The language of $T_1$ and $T_2$ relaxation, familiar from NMR
spectroscopy, clarifies what can and cannot be captured by TDDFT.
Consider a two-level system with ground state $|g\rangle$ and excited
state $|e\rangle$.  The density matrix is:
\begin{equation}
  \rho = \begin{pmatrix} \rho_{ee} & \rho_{eg} \\ \rho_{ge} & \rho_{gg} \end{pmatrix},
  \qquad \rho_{ee} + \rho_{gg} = 1.
\end{equation}
The Bloch equations with relaxation read:
\begin{align}
  \dot{\rho}_{ee} &= -\frac{\rho_{ee}}{T_1}
    + \text{(driving terms)}, \label{eq:bloch_pop}\\
  \dot{\rho}_{eg} &= -\ii\omega_0\rho_{eg}
    - \frac{\rho_{eg}}{T_2}
    + \text{(driving terms)}, \label{eq:bloch_coh}
\end{align}
where $\omega_0 = E_e - E_g$ is the transition frequency.

\begin{itemize}
  \item $T_1$ (\textbf{longitudinal, population relaxation}): excited
    population $\rho_{ee}$ decays exponentially.  In spectroscopy this
    is spontaneous emission + non-radiative decay.  Physical mechanism:
    coupling to the radiation field (QED photon bath) or phonons.
    \emph{Not present in TDDFT, which describes only electrons.}

  \item $T_2$ (\textbf{transverse, dephasing}): the off-diagonal element
    $\rho_{eg}$ (coherence) decays as $e^{-t/T_2}$.  The general bound
    is $T_2 \leq 2T_1$.  The excess decay beyond the $2T_1$ limit is
    called \emph{pure dephasing} ($T_2^*$) and arises from
    elastic processes that randomize the quantum phase without
    changing populations:
    \begin{equation}
      \frac{1}{T_2} = \frac{1}{2T_1} + \frac{1}{T_2^*}.
      \label{eq:T2}
    \end{equation}
\end{itemize}

\begin{insight}
\textbf{Your atom-in-vacuum example is exactly right.}
A single atom in perfect vacuum with no external coupling has:
$T_1 = \infty$ (no radiation bath to emit into --- spontaneous emission
is a QED effect, absent from the Schr\"{o}dinger equation for electrons
alone), and $T_2 = \infty$ (no elastic collisions to dephase the
coherence).  The excited state is stable forever.  This is the
correct quantum-mechanical result in the absence of a bath.
\end{insight}

\subsection{The Reduced Density Matrix and Intrinsic Dephasing}

Here is the resolution of the apparent paradox between unitary evolution
and dephasing.

The exact $N$-electron wavefunction $\Psi(\bm{r}_1,\ldots,\bm{r}_N,t)$
evolves unitarily.  The \emph{one-body reduced density matrix} (1-RDM):
\begin{equation}
  \gamma(\bm{r},\bm{r}',t) = N\int\!\Psi^*(\bm{r},\bm{r}_2,\ldots)\,
  \Psi(\bm{r}',\bm{r}_2,\ldots)\,\dd\bm{r}_2\cdots\dd\bm{r}_N
  \label{eq:1rdm}
\end{equation}
is obtained by tracing out $N-1$ electrons.  Its equation of motion is:
\begin{equation}
  \ii\dot{\gamma} = \comm{T + V_{\mathrm{ext}}}{\gamma} + \mathcal{C},
  \label{eq:1rdm_eom}
\end{equation}
where $\mathcal{C}$ is the \emph{collision integral} --- the contribution
from the two-body correlations that survive after tracing over $N-1$
electrons.  This term is \emph{non-unitary} for the 1-RDM even though
the $N$-body evolution is unitary.  It is the source of intrinsic
dephasing from electron-electron scattering.

\textbf{Contrast with the KS density matrix.}
The KS system is \emph{non-interacting} by construction: its $N$ electrons
move independently in the effective potential $V_\text{KS}$.
Tracing over $N-1$ KS electrons produces no residual two-body correlation
because there are none --- so $\mathcal{C}=0$ for the KS system.
The KS density matrix therefore satisfies the \emph{homogeneous} equation:
\begin{equation}
  \ii\dot{P} = \comm{\Hks}{P}, \qquad \mathcal{C} = 0.
\end{equation}
Both $\gamma$ and $P$ satisfy equations that \emph{look} like $i\dot{M}=[H,M]$,
but they describe different objects at different levels of theory.
The collision integral $\mathcal{C}$ is precisely what $P$ is missing
compared to $\gamma$, and it is what would be needed for physical dephasing.

\begin{derivbox}
\textbf{Why $\mathcal{C}$ causes pure dephasing.}

Decompose $\mathcal{C}$ into a Hermitian part $\mathcal{C}_H$ and an
anti-Hermitian part $\mathcal{C}_A$:
\[
  \mathcal{C} = \mathcal{C}_H + \mathcal{C}_A,
  \quad
  \mathcal{C}_H = \tfrac{1}{2}(\mathcal{C}+\mathcal{C}^\dagger),
  \quad
  \mathcal{C}_A = \tfrac{1}{2}(\mathcal{C}-\mathcal{C}^\dagger).
\]
$\mathcal{C}_H$ renormalizes the effective Hamiltonian (energy shifts).
$\mathcal{C}_A$ is non-Hermitian and causes the \emph{magnitude} of
off-diagonal elements to decay:
\[
  \ii\dot{\gamma}_{ij} = (\epsilon_i - \epsilon_j)\gamma_{ij}
  + \mathcal{C}_{A,ij}.
\]
If $\mathcal{C}_A = -\ii\Gamma\,\gamma$ with $\Gamma > 0$, then
$\gamma_{ij}(t) \propto e^{-\ii\omega_{ij}t}e^{-\Gamma t}$ ---
the coherence decays exponentially while the diagonal populations
(in the eigenstate basis) remain conserved.  This is pure dephasing
with $T_2^* = 1/\Gamma$.
\end{derivbox}

In TDDFT, the exact $\fxc(\omega)$ with
$\operatorname{Im}\fxc(\omega)\neq 0$ plays the role of $\mathcal{C}_A$.
Here is why.
The linearised TDKS equation for a density perturbation $\delta\rho$ reads
(in frequency space, using Eq.~\ref{eq:func_taylor}):
\begin{equation}
  \delta\rho(\omega)
  = \chi_s\!\left[\delta V_{\mathrm{ext}}
    + \frac{1}{|\bm{r}-\bm{r}'|}\delta\rho
    + \fxc(\omega)\,\delta\rho\right]\!(\omega).
\end{equation}
When $\fxc(\omega)$ is real, the effective KS Hamiltonian that drives
$\delta\rho$ is Hermitian, and the response is purely oscillatory (no
decay).
When $\operatorname{Im}\fxc(\omega)\neq 0$, the effective driving potential
acquires an imaginary part, which acts like a damping term in the equation
of motion for the off-diagonal elements of $P$.  Specifically, an
off-diagonal coherence $P_{nm}$ evolves as
$\ii\dot{P}_{nm} \approx (\omega_{nm} - \ii\,\Gamma_{nm})P_{nm}$,
where $\Gamma_{nm} \propto \operatorname{Im}\fxc(\omega_{nm}) > 0$
is the dephasing rate.  This is the same exponential decay produced by
the collision integral $\mathcal{C}_A$ in the reduced density matrix
equation:
it introduces an effective non-Hermitian correction to the KS
Hamiltonian in the frequency domain, providing $T_2$ dephasing without
violating $N$-body unitarity.

\paragraph{Summary.}
\begin{center}
\begin{tabular}{@{}lll@{}}
\toprule
Relaxation & Physical origin & In TDDFT? \\
\midrule
$T_1$ (population) & Coupling to radiation/phonon bath & No \\
$T_2$ pure dephasing & Electron-electron scattering ($\mathcal{C}_A$) & Only with $\operatorname{Im}\fxc\neq 0$ \\
$T_2$ from $T_1$ & Same bath as $T_1$ & No \\
\bottomrule
\end{tabular}
\end{center}

% ============================================================================
\section{Double Excitations: Selection Rules and CI Mixing}
\label{sec:doubles}
% ============================================================================

\subsection{Pure Double Excitations Are Dipole-Forbidden --- Your Intuition Is Correct}

Consider two non-interacting helium atoms A and B, separated in space.
Define:
\begin{align*}
  |G_A\rangle &= \text{He}_A \text{ ground state } (1s)^2, \\
  |E_A\rangle &= \text{He}_A \text{ singly excited state } (1s)(2p), \\
  |G_B\rangle,\;|E_B\rangle &= \text{same for He}_B.
\end{align*}
The composite two-atom states are tensor products:
\begin{align*}
  |\Psi_0\rangle &= |G_A\rangle\otimes|G_B\rangle \equiv |GG\rangle, \\
  |\Psi_{\rm S}^A\rangle &= |E_A\rangle\otimes|G_B\rangle \equiv |EG\rangle, \\
  |\Psi_{\rm S}^B\rangle &= |G_A\rangle\otimes|E_B\rangle \equiv |GE\rangle, \\
  |\Psi_{\rm D}\rangle &= |E_A\rangle\otimes|E_B\rangle \equiv |EE\rangle.
\end{align*}
The dipole operator for the whole system decomposes as:
\begin{equation}
  \hat{x} = \hat{x}_A\otimes\hat{1}_B + \hat{1}_A\otimes\hat{x}_B.
  \label{eq:dipole_total}
\end{equation}

\begin{derivbox}
\textbf{Transition moment from ground to doubly excited state.}
\begin{align}
  \mel{EE}{\hat{x}}{GG}
  &= \mel{EE}{\hat{x}_A\otimes\hat{1}_B + \hat{1}_A\otimes\hat{x}_B}{GG}
  \notag\\
  &= \underbrace{\mel{E_A}{\hat{x}_A}{G_A}}_{\neq 0}
     \underbrace{\braket{E_B}{G_B}}_{=0}
   + \underbrace{\braket{E_A}{G_A}}_{=0}
     \underbrace{\mel{E_B}{\hat{x}_B}{G_B}}_{\neq 0}
  \notag\\
  &= 0.
  \label{eq:double_forbidden}
\end{align}
Both terms vanish because excited and ground states on the
\emph{other} atom are orthogonal.  The dipole operator connects only
states that differ by a \emph{single} local excitation; reaching $|EE\rangle$
requires both atoms to be excited simultaneously, which a one-body
operator cannot achieve.
\end{derivbox}

\begin{insight}
You are correct: a pure doubly excited state is dipole-dark from the
ground state.  To reach $|EE\rangle$ directly requires a two-photon
operator $\hat{x}_A\hat{x}_B$ --- which is accessible in two-photon
absorption spectroscopy.  This is a different (and real) phenomenon.

The TDDFT ``double excitation problem'' concerns a \emph{different}
situation: states that are \textbf{predominantly} doubly excited but
have \textbf{mixed} character due to Coulomb coupling, and that
\emph{are} dipole-accessible.
\end{insight}

\subsection{Turning On the Coulomb Interaction: CI Mixing}

Now allow a Coulomb interaction between the two He atoms:
\begin{equation}
  V_{ee}^{AB} = \sum_{i\in A}\sum_{j\in B} \frac{1}{|\bm{r}_i - \bm{r}_j|}.
  \label{eq:vee}
\end{equation}
Treat $V_{ee}^{AB}$ as a perturbation.  By first-order degenerate
perturbation theory, the eigenstates of $H_0 + V_{ee}^{AB}$ are
linear combinations of the unperturbed states.  In the subspace
$\{|EG\rangle, |GE\rangle, |EE\rangle, \ldots\}$:

\begin{derivbox}
\textbf{First-order CI mixing of $|EE\rangle$ into $|EG\rangle$.}

The coupling matrix element between $|EG\rangle$ and $|EE\rangle$
is:
\[
  \mel{EG}{V_{ee}^{AB}}{EE}
  = \mel{E_A,G_B}{V_{ee}^{AB}}{E_A,E_B}.
\]
Since $V_{ee}^{AB}$ involves coordinates of both atoms, this is
generically nonzero --- it is the inter-atomic Coulomb integral.

To first order in $V_{ee}^{AB}$, the corrected eigenstate is:
\begin{equation}
  |\tilde{\Psi}_1\rangle
  = \alpha|EG\rangle + \beta|GE\rangle
  + \underbrace{\frac{\mel{EE}{V_{ee}^{AB}}{EG}}{E_{EG}-E_{EE}}}_{\gamma}|EE\rangle
  + \cdots
  \label{eq:mixed_state}
\end{equation}
The coefficient $\gamma$ is generically nonzero; the state
$|\tilde{\Psi}_1\rangle$ has genuine double-excitation character.
\end{derivbox}

Now compute the dipole transition moment from the ground state to
$|\tilde{\Psi}_1\rangle$:
\begin{align}
  \mel{\tilde{\Psi}_1}{\hat{x}}{GG}
  &= \alpha\underbrace{\mel{EG}{\hat{x}}{GG}}_{\neq 0}
   + \beta\underbrace{\mel{GE}{\hat{x}}{GG}}_{\neq 0}
   + \gamma\underbrace{\mel{EE}{\hat{x}}{GG}}_{=0}
  \notag\\
  &= \alpha\mel{E_A}{\hat{x}_A}{G_A}
   + \beta\mel{E_B}{\hat{x}_B}{G_B}
  \neq 0.
  \label{eq:mixed_dipole}
\end{align}

\begin{keyresult}
Once the Coulomb interaction mixes a double-excitation component
into an otherwise singly-excited state, that state becomes:
\begin{enumerate}[noitemsep]
  \item \textbf{Dipole-accessible} from the ground state (nonzero
    transition moment), and
  \item \textbf{At an energy shifted} by the CI mixing --- an energy
    that adiabatic TDDFT cannot reproduce.
\end{enumerate}
The TDDFT failure is primarily \textbf{energetic}: the state appears
in the spectrum, but at the wrong position.
\end{keyresult}

\subsection{The Casida Linear-Response Equations}
\label{sec:casida}

Linear-response TDDFT (Casida 1995) reframes excitation energies as
eigenvalues of a matrix equation.  Starting from the linear response
of the TDKS density to an external perturbation $V_{\rm ext}(\omega)$,
one derives the response function:
\begin{equation}
  \delta\rho(\omega) = \chi(\omega)\,V_{\rm ext}(\omega),
  \qquad
  \chi(\omega) = \chi_0(\omega)\bigl[1 - (v_{ee}+\fxc(\omega))\chi_0(\omega)\bigr]^{-1},
  \label{eq:chi}
\end{equation}
where $\chi_0$ is the non-interacting KS response function and
$v_{ee}(\bm{r},\bm{r}') = 1/|\bm{r}-\bm{r}'|$.

Poles of $\chi(\omega)$ are the excitation frequencies.  In the
molecular orbital (MO) basis $\{\phi_i, \phi_a\}$ (occupied $i$,
virtual $a$), expanding $\chi_0$ in terms of KS orbital pairs gives
the Casida eigenvalue problem:
\begin{equation}
  \begin{pmatrix} A(\omega) & B \\ B^* & A^*(\omega) \end{pmatrix}
  \begin{pmatrix} X \\ Y \end{pmatrix}
  = \omega\begin{pmatrix} \mathbf{1} & 0 \\ 0 & -\mathbf{1} \end{pmatrix}
  \begin{pmatrix} X \\ Y \end{pmatrix},
  \label{eq:casida}
\end{equation}
where:
\begin{align}
  A_{ia,jb}(\omega) &= (\epsilon_a-\epsilon_i)\delta_{ij}\delta_{ab}
    + K_{ia,jb}(\omega),\label{eq:Amat}\\
  B_{ia,jb} &= K_{ia,bj},\\
  K_{ia,jb}(\omega) &= \iint\phi_i^*(\bm{r})\phi_a(\bm{r})
    \left[\frac{1}{|\bm{r}-\bm{r}'|}+\fxc(\bm{r},\bm{r}',\omega)\right]
    \phi_j(\bm{r}')\phi_b^*(\bm{r}')\,\dd^3r\,\dd^3r'.
\end{align}

\begin{insight}
\textbf{Connection to Gaussian-basis CI/RPA: $K_{ia,jb}$ and the ERI $(ia|jb)$.}
The Coulomb part of $K_{ia,jb}$ is exactly a two-electron repulsion integral
(ERI) in the MO basis:
\begin{equation}
  K_{ia,jb}^{\rm Coulomb}
  = \iint\frac{\phi_i^*(\bm{r})\phi_a(\bm{r})\,
               \phi_j(\bm{r}')\phi_b^*(\bm{r}')}{|\bm{r}-\bm{r}'|}
    \,\dd^3r\,\dd^3r'
  \equiv (ia|jb).
  \label{eq:Kij_ERI}
\end{equation}
In Gaussian-basis codes this integral is stored in the two-electron integral
file and is familiar from CIS, RPA, and TDDFT calculations.
The full Casida matrix is therefore:
\begin{equation}
  K_{ia,jb}(\omega)
  = (ia|jb) + (ia|\fxc|jb),
  \label{eq:K_split}
\end{equation}
where $(ia|\fxc|jb) = \iint\phi_i^*\phi_a\,\fxc\,\phi_j\phi_b^*\,\dd^3r\,\dd^3r'$
is the XC correction.

\begin{center}
\renewcommand{\arraystretch}{1.3}
\begin{tabular}{@{}lll@{}}
\toprule
Method & $K_{ia,jb}$ & Notes \\
\midrule
CIS (HF) & $(ia|jb) - (ij|ab)$ & Coulomb + exchange, no correlation \\
RPA (adiabatic) & $(ia|jb)$ & Coulomb only (Tamm-Dancoff approx.) \\
TDDFT-LDA & $(ia|jb) + (ia|f_{\rm xc}^{\rm LDA}|jb)$ & Coulomb + local XC \\
Exact TDDFT & $(ia|jb) + (ia|\fxc(\omega)|jb)$ & frequency-dep.\ XC kernel \\
\bottomrule
\end{tabular}
\renewcommand{\arraystretch}{1}
\end{center}
The real-space grid in RMG evaluates the same integrals numerically:
$\sum_i \phi_i^*(\bm{r}_i)\phi_a(\bm{r}_i)\,V_H(\bm{r}_i)\,h^3$
is the grid-point analogue of $(ia|jb)$ contracted over the Hartree
potential.
\end{insight}

\paragraph{Dimension of the Casida matrix.}
The indices $(i,a)$ and $(j,b)$ run over \emph{single} KS
excitation pairs: $N_{\rm occ}\times N_{\rm virt}$ in total.
\textbf{Double excitations $(i,j)\to(a,b)$ are not included as
basis vectors.}

\subsection{Why Double Excitations Require Poles in $\fxc(\omega)$}

To see why a pole in $\fxc$ is needed, consider what happens when a
doubly excited state $|D\rangle$ at energy $\Omega_D$ is near a
singly excited KS pair $(i\to a)$ at energy $\omega_{ia} = \epsilon_a-\epsilon_i$.

\begin{derivbox}
\textbf{A pole in $\fxc$ is equivalent to including $|D\rangle$ in CI.}

Maitra \etal~(2004) showed that the contribution of the doubly
excited state to the exact $\fxc$ takes the form:
\begin{equation}
  \fxc(\omega) \sim \frac{|\mel{\Phi_{ia}}{V_{ee}}{\Phi_{D}}\,|^2}
                        {\omega - \omega_D^{\rm KS}}
  + \text{regular},
  \label{eq:maitra_pole}
\end{equation}
where $\Phi_{ia}$ is the KS single-excitation Slater determinant and
$\Phi_D$ is the doubly excited determinant.  The numerator is the
Coulomb matrix element coupling single to double excitation --- exactly
the CI matrix element.

Substituting Eq.~(\ref{eq:maitra_pole}) into the Casida coupling matrix
$K_{ia,jb}(\omega)$ at $\omega\approx\omega_D^{\rm KS}$ gives a
divergence that ``pulls'' a new eigenvalue of Eq.~(\ref{eq:casida})
to the position $\Omega_D$ --- reproducing the double excitation.

With the \emph{adiabatic} kernel $\fxc(\omega) = \text{const}$, the
pole is absent, and $\Omega_D$ never appears as an eigenvalue,
regardless of the strength of the driving field.
\end{derivbox}

\begin{remark}
Using a field with symmetry $\hat{x}^2$ or $\hat{x}\hat{y}$ (two-photon
absorption) drives a different physical process: it accesses states
that are two-photon-allowed, which is a distinct spectroscopic
phenomenon from what the Casida equations describe.
The ``missing double excitations'' in TDDFT are states accessible by
\emph{one photon} whose energy is wrong because of missing CI mixing.
\end{remark}

\paragraph{Energy shift from CI: perturbation theory estimate.}
The second-order energy correction from the coupling between the
singly excited state $|S\rangle$ at $\omega_S$ and the doubly excited
state $|D\rangle$ at $\omega_D$ is:
\begin{equation}
  \Delta\omega_S = \frac{|\mel{S}{V_{ee}}{D}|^2}{\omega_S - \omega_D}.
  \label{eq:2nd_order}
\end{equation}
Adiabatic TDDFT misses this shift entirely.  The magnitude can be
$0.2$--$0.5$\,eV in conjugated molecules, explaining why the
TDDFT absorption spectrum of polyenes is systematically too high in energy.

% ============================================================================
\section{Gauge Invariance and Current-Density TDDFT}
\label{sec:gauge}
% ============================================================================

\subsection{Gauge Transformations in Quantum Mechanics}

The Hamiltonian for an electron in an electromagnetic field is
(in Gaussian units converted to atomic units):
\begin{equation}
  H = \frac{1}{2}\bigl(\hat{\bm{p}} + \bm{A}(\bm{r},t)\bigr)^2
  + \phi(\bm{r},t),
  \label{eq:H_em}
\end{equation}
where $\bm{A}$ is the vector potential and $\phi$ is the scalar potential.
The physical electric and magnetic fields are:
\begin{equation}
  \bm{E} = -\nabla\phi - \frac{\partial\bm{A}}{\partial t},
  \qquad
  \bm{B} = \nabla\times\bm{A}.
  \label{eq:EB}
\end{equation}
A \emph{gauge transformation} with a scalar function $\chi(\bm{r},t)$ leaves
$\bm{E}$ and $\bm{B}$ unchanged while transforming the potentials:
\begin{align}
  \bm{A} &\to \bm{A}' = \bm{A} + \nabla\chi, \label{eq:gauge_A}\\
  \phi &\to \phi' = \phi - \frac{\partial\chi}{\partial t}. \label{eq:gauge_phi}
\end{align}
The wavefunction transforms as:
\begin{equation}
  \Psi(\bm{r},t) \to \Psi'(\bm{r},t) = e^{\ii\chi(\bm{r},t)}\,\Psi(\bm{r},t).
  \label{eq:gauge_psi}
\end{equation}
One can verify directly that $\Psi'$ satisfies the Schr\"{o}dinger equation with
$(H', \Psi')$ if and only if $\Psi$ satisfies it with $(H, \Psi)$.

\subsection{The Charge Density Is Gauge-Invariant}

\begin{derivbox}
Under Eq.~(\ref{eq:gauge_psi}):
\begin{equation}
  \rho'(\bm{r},t) = |\Psi'|^2
  = \bigl|e^{\ii\chi}\Psi\bigr|^2
  = e^{-\ii\chi}e^{+\ii\chi}|\Psi|^2
  = |\Psi|^2 = \rho(\bm{r},t).
  \label{eq:rho_invariant}
\end{equation}
The charge density is \textbf{gauge-invariant}: two different gauges
related by $\chi$ give identical $\rho(\bm{r},t)$ for all $t$.
\end{derivbox}

\subsection{The Paramagnetic Current Is NOT Gauge-Invariant}

The paramagnetic (canonical) current density is defined as:
\begin{equation}
  \bm{J}_p(\bm{r},t) = \frac{1}{2\ii}\bigl[
    \Psi^*(\bm{r},t)\,\nabla\Psi(\bm{r},t)
    - \Psi(\bm{r},t)\,\nabla\Psi^*(\bm{r},t)
  \bigr].
  \label{eq:Jp}
\end{equation}

\begin{derivbox}
Under the gauge transformation $\Psi\to e^{\ii\chi}\Psi$, compute
$\nabla\Psi'$:
\begin{equation}
  \nabla\Psi' = \nabla(e^{\ii\chi}\Psi)
  = \ii(\nabla\chi)\,e^{\ii\chi}\Psi + e^{\ii\chi}\nabla\Psi.
  \label{eq:grad_psi_gauge}
\end{equation}
Substituting into Eq.~(\ref{eq:Jp}):
\begin{align}
  \bm{J}_p' &= \frac{1}{2\ii}\bigl[
    (e^{-\ii\chi}\Psi^*)
    \bigl(\ii(\nabla\chi)e^{\ii\chi}\Psi + e^{\ii\chi}\nabla\Psi\bigr)
    - \text{c.c.}
  \bigr] \notag\\
  &= \frac{1}{2\ii}\bigl[
    \Psi^*\nabla\Psi - \Psi\nabla\Psi^*
  \bigr]
  + \frac{1}{2\ii}\bigl[
    \ii\nabla\chi|\Psi|^2 + \ii\nabla\chi|\Psi|^2
  \bigr] \notag\\
  &= \bm{J}_p + \rho\,\nabla\chi.
  \label{eq:Jp_gauge}
\end{align}
\end{derivbox}

\begin{keyresult}
$\bm{J}_p$ is \textbf{not} gauge-invariant: it changes by
$\rho\,\nabla\chi$ under the transformation.
\end{keyresult}

\subsection{The Physical (Total) Current Is Gauge-Invariant}

\paragraph{Canonical vs kinetic momentum.}
In classical mechanics, a charged particle in a vector potential $\bm{A}$
has two distinct momenta:
\begin{itemize}[noitemsep]
  \item \textbf{Canonical momentum}: $\bm{p} = m\bm{v} + e\bm{A}$
    (the quantity conjugate to position in the Lagrangian; gauge-dependent).
  \item \textbf{Kinetic momentum}: $m\bm{v} = \bm{p} - e\bm{A}$
    (the physically observable momentum that equals $m\times$velocity;
    gauge-invariant).
\end{itemize}
The same distinction carries over to quantum mechanics.
The paramagnetic current $\bm{J}_p$ (Section~\ref{sec:Jp_not_invariant})
is built from $\bm{p} = -i\nabla$ — the \emph{canonical} momentum operator.
The physical, measurable current is built from the \emph{kinetic} momentum
$(-i\nabla - e\bm{A})$, giving:
\begin{equation}
  \bm{J}(\bm{r},t)
  = \underbrace{\bm{J}_p(\bm{r},t)}_{\text{canonical part}}
  - \underbrace{\rho(\bm{r},t)\,\bm{A}(\bm{r},t)}_{\text{diamagnetic part}}.
  \label{eq:J_total}
\end{equation}
The second term $-\rho\bm{A}$ is called the \textbf{diamagnetic current}:
it represents the response of the charge density to the vector potential
itself, and it is what makes $\bm{J}$ gauge-invariant while $\bm{J}_p$
alone is not.

\begin{derivbox}
Under the gauge transformation $\bm{A}\to\bm{A}+\nabla\chi$,
$\bm{J}_p\to\bm{J}_p + \rho\nabla\chi$:
\begin{align}
  \bm{J}' &= \bm{J}_p' - \rho'\bm{A}' \notag\\
  &= (\bm{J}_p + \rho\nabla\chi) - \rho(\bm{A}+\nabla\chi) \notag\\
  &= \bm{J}_p - \rho\bm{A} = \bm{J}.
  \label{eq:J_invariant}
\end{align}
\end{derivbox}

\begin{keyresult}
$\bm{J}(\bm{r},t) = \bm{J}_p - \rho\bm{A}$ is \textbf{gauge-invariant}.
It is the physically observable current that appears in Maxwell's equations.
\end{keyresult}

\subsection{Why $\rho$ Alone Fails for Periodic Systems}

For an isolated finite system, the Runge-Gross theorem establishes a
one-to-one mapping $V_{\rm ext}(t)\leftrightarrow\rho(t)$ even without
discussing the gauge.  The reason: for a finite system, one can always
choose a unique gauge by requiring, e.g., $\bm{A}=0$ (the Coulomb gauge).

For a \emph{periodic crystal}, a spatially uniform (but time-dependent)
vector potential $\bm{A}(t)$ drives a macroscopic current.  Under a gauge
transformation with $\chi(\bm{r},t) = -\bm{A}(t)\cdot\bm{r}$:
\begin{align}
  \bm{A}(t) &\to \bm{A}' = \bm{A}(t) + \nabla(-\bm{A}(t)\cdot\bm{r})
              = \bm{A}(t) - \bm{A}(t) = 0, \\
  \phi &\to \phi' = \phi + \dot{\bm{A}}(t)\cdot\bm{r} = \bm{E}(t)\cdot\bm{r}.
\end{align}
This gauge transformation relates the \emph{velocity gauge}
(vector potential $\bm{A}$, no scalar potential) to the \emph{length gauge}
(scalar potential $\bm{E}\cdot\bm{r}$, no vector potential).

The critical observation: in a crystal, $\chi = -\bm{A}(t)\cdot\bm{r}$ is
\emph{spatially linear}, so $\nabla\chi$ is uniform.
A uniform $\nabla\chi$ does not break the periodicity of $\rho(\bm{r},t)$:
the density remains periodic regardless of the gauge.
Therefore, \textbf{measuring $\rho(\bm{r},t)$ alone cannot distinguish the
velocity gauge from the length gauge} --- it cannot detect whether $\bm{A}(t)$
is present.  But $\bm{J}(\bm{r},t)$ can: the macroscopic current changes
sign when $\bm{A}(t)$ reverses.

This is the fundamental reason why the Runge-Gross theorem in its original
form does not apply to periodic solids: the mapping
$\bm{A}(t)\leftrightarrow\rho(t)$ is not one-to-one.

\subsection{The Vignale-Kohn Theorem}

Vignale and Kohn (1996) extended the Runge-Gross argument to the pair
$(\rho(\bm{r},t),\,\bm{J}(\bm{r},t))$:

\begin{theorem}[Vignale-Kohn 1996]
Given a fixed initial state $\Psi_0$, there is a one-to-one mapping
between the physical electromagnetic field and the pair
$(\rho(\bm{r},t),\,\bm{J}(\bm{r},t))$.  Both $\rho$ and $\bm{J}$ together
uniquely determine all physical observables.
\end{theorem}

Since both $\rho$ and $\bm{J}$ are gauge-invariant (proved above), this
theorem is gauge-consistent.  The basic variable of current-density TDDFT
(CDFT) is therefore the pair $(\rho, \bm{J})$, and the XC functional
is:
\begin{equation}
  \Vxc^{\rm CDFT}[\rho,\bm{J};\,\bm{r},t]
  = \Vxc^{\rm CDFT}[\rho(\bm{r}',t'),\bm{J}(\bm{r}',t');\,t'\leq t].
  \label{eq:vxc_cdft}
\end{equation}

% -----------------------------------------------------------------------
\subsection{What Is a Kernel?}
\label{sec:kernel}
% -----------------------------------------------------------------------

The word \emph{kernel} appears repeatedly in this document --- the XC
kernel, the memory kernel, the viscosity kernel.  Before applying these
concepts to the electron gas, we build the mathematical intuition from
the ground up.

\subsubsection{Integral Operators and Their Kernels}

An \emph{integral operator} $\mathcal{K}$ maps an input function $f$
to an output function $g$ via integration:
\begin{equation}
  g(x) = (\mathcal{K}f)(x)
  = \int K(x,\,y)\,f(y)\,\dd y.
  \label{eq:integral_op}
\end{equation}
The function $K(x,y)$ is the \emph{kernel} of the operator.  It encodes
the rule: ``how much does the value of $f$ at point $y$ contribute to
the result at point $x$?''

\paragraph{Three everyday examples.}
\begin{enumerate}
  \item \textbf{Gaussian smoothing} (image processing):
    \begin{equation}
      g(x) = \int \frac{1}{\sqrt{2\pi}\sigma}
      e^{-(x-y)^2/2\sigma^2}\,f(y)\,\dd y.
      \label{eq:gaussian_kernel}
    \end{equation}
    The kernel $K(x,y) = \frac{1}{\sqrt{2\pi}\sigma}e^{-(x-y)^2/2\sigma^2}$
    is a Gaussian centered at $y=x$.  Each output value is a weighted
    average of nearby input values.  Points far from $x$ contribute
    little; points close to $x$ contribute much.

  \item \textbf{Hartree (Coulomb) potential}:
    \begin{equation}
      V_H(\bm{r}) = \int \frac{\rho(\bm{r}')}{|\bm{r}-\bm{r}'|}\,\dd^3r'.
      \label{eq:hartree_kernel}
    \end{equation}
    The kernel is $K(\bm{r},\bm{r}') = 1/|\bm{r}-\bm{r}'|$.
    The Hartree potential at $\bm{r}$ is the sum of Coulomb contributions
    from the density at every other point $\bm{r}'$.  This is a
    \emph{nonlocal} operator in space: $V_H(\bm{r})$ depends on $\rho$
    everywhere, not just at $\bm{r}$.

  \item \textbf{Green's function of the Laplacian}:
    The Poisson equation $-\nabla^2 V = 4\pi\rho$ has the solution
    $V(\bm{r}) = \int G(\bm{r},\bm{r}')\rho(\bm{r}')\,\dd^3r'$
    with $G(\bm{r},\bm{r}') = 1/|\bm{r}-\bm{r}'|$.
    Green's functions are kernels of the inverse differential operator.
\end{enumerate}

\subsubsection{Convolution Kernels: Translation-Invariant Operators}

When the kernel depends only on the difference $K(x,y)=K(x-y)$, the
integral operator is a \emph{convolution}:
\begin{equation}
  g(x) = \int K(x-y)\,f(y)\,\dd y \equiv (K*f)(x).
  \label{eq:convolution}
\end{equation}
Convolution kernels are particularly natural for:
\begin{itemize}[noitemsep]
  \item Homogeneous systems (same physics at every point).
  \item Time-domain response (system looks the same at every time $t$).
  \item Frequency analysis: convolution in time $\leftrightarrow$
    \emph{multiplication} in frequency, $\hat{g}(\omega) = \hat{K}(\omega)\hat{f}(\omega)$.
\end{itemize}

\subsubsection{Causal (Memory) Kernels}

In physics, response functions must be \emph{causal}: the effect cannot
precede its cause.  A causal convolution kernel satisfies:
\begin{equation}
  K(\tau) = 0 \quad\text{for } \tau < 0,
  \label{eq:causal}
\end{equation}
so the response at time $t$ depends only on inputs at times $t'<t$:
\begin{equation}
  g(t) = \int_0^t K(t-t')\,f(t')\,\dd t'
  = \int_0^\infty K(\tau)\,f(t-\tau)\,\dd\tau.
  \label{eq:causal_conv}
\end{equation}
The kernel $K(\tau)$ tells you ``how strongly does an input from $\tau$
seconds ago influence the current output?''  If $K(\tau)\propto e^{-\tau/\tau_0}$,
the influence decays exponentially with a characteristic memory time
$\tau_0$.

\begin{insight}
A system with a causal memory kernel is the mathematical formalization
of your car analogy.  The kernel $K(\tau)$ weights the past trajectory
$f(t')$ of the input --- car position, density, deformation --- and the
current output $g(t)$ depends on the full weighted history.  The
adiabatic approximation corresponds to $K(\tau)=c\,\delta(\tau)$: zero
memory, only the instantaneous value matters.
\end{insight}

\subsubsection{The XC Kernel in DFT}

The \emph{exchange-correlation kernel} $\fxc$ is the functional
derivative of the XC potential with respect to the density:
\begin{equation}
  \fxc(\bm{r},\bm{r}',t-t')
  = \frac{\delta\Vxc(\bm{r},t)}{\delta\rho(\bm{r}',t')}.
  \label{eq:fxc_kernel}
\end{equation}
This answers the question: ``if I add a tiny density perturbation
$\delta\rho(\bm{r}')$ at time $t'$, how does $\Vxc(\bm{r})$ respond at
a later time $t$?''

In the \emph{static} (ground-state) limit, the XC kernel reduces to
the second functional derivative of the XC energy:
\begin{equation}
  \fxc^{\rm gs}(\bm{r},\bm{r}')
  = \frac{\delta^2 \Exc[\rho]}{\delta\rho(\bm{r})\,\delta\rho(\bm{r}')}.
  \label{eq:fxc_static_def}
\end{equation}
For LDA, where $\Exc = \int\epsilon_{\rm xc}(n)\rho\,\dd^3r$:
\begin{equation}
  \fxc^{\rm LDA}(\bm{r},\bm{r}')
  = \frac{\dd^2[\rho\,\epsilon_{\rm xc}(\rho)]}{\dd\rho^2}\,\delta(\bm{r}-\bm{r}').
  \label{eq:fxc_lda}
\end{equation}
The delta function $\delta(\bm{r}-\bm{r}')$ makes LDA \emph{local} in space:
the XC potential at $\bm{r}$ depends only on the density at the same
point $\bm{r}$.  GGA kernels have a spatial gradient correction
(nonzero for $|\bm{r}-\bm{r}'|\lesssim$ grid spacing).
The frequency-dependent kernel $\fxc(\omega)$ additionally captures
nonlocality \emph{in time} --- the memory.

% -----------------------------------------------------------------------
\subsection{Elements of Continuum Mechanics}
\label{sec:mechanics}
% -----------------------------------------------------------------------

The Vignale-Kohn functional treats the electron gas as a
\emph{viscoelastic fluid} --- a material that has both elastic
(spring-like) and viscous (dashpot-like) properties, with a
memory of its deformation history.  This section builds the
necessary continuum mechanics background from scratch.

\subsubsection{Displacement Field and Velocity Field}

Consider a fluid (or solid) in which each material element
at position $\bm{r}_0$ (its reference position) is displaced to
a new position $\bm{r}_0 + \bm{u}(\bm{r}_0,t)$ at time $t$.
The vector field $\bm{u}(\bm{r},t)$ is the \emph{displacement field}.

The \emph{velocity field} of the fluid elements is:
\begin{equation}
  \bm{v}(\bm{r},t) = \frac{\partial\bm{u}(\bm{r},t)}{\partial t}.
  \label{eq:velocity}
\end{equation}
For the electron gas, $\bm{u}(\bm{r},t)$ is the displacement of the
electron fluid element at $\bm{r}$, and the current density is:
\begin{equation}
  \bm{J}(\bm{r},t) = n_0(\bm{r})\,\bm{v}(\bm{r},t)
  = n_0(\bm{r})\,\frac{\partial\bm{u}}{\partial t},
  \label{eq:current_displacement}
\end{equation}
where $n_0(\bm{r})$ is the equilibrium electron density.

\subsubsection{The Strain Tensor: Measuring Deformation}

The \emph{strain tensor} $\varepsilon_{\alpha\beta}$ measures how much
the material is locally deformed:
\begin{equation}
  \varepsilon_{\alpha\beta} = \frac{1}{2}
  \left(\frac{\partial u_\alpha}{\partial r_\beta}
      + \frac{\partial u_\beta}{\partial r_\alpha}\right).
  \label{eq:strain}
\end{equation}
It is symmetric ($\varepsilon_{\alpha\beta} = \varepsilon_{\beta\alpha}$)
and captures two types of deformation:

\paragraph{Volumetric (compressive) part --- the trace.}
The trace of the strain tensor measures volume change:
\begin{equation}
  \varepsilon_{\rm vol} = \operatorname{Tr}(\varepsilon)
  = \nabla\cdot\bm{u}
  = \frac{\delta V}{V_0}.
  \label{eq:volumetric}
\end{equation}
$\nabla\cdot\bm{u} > 0$: expansion; $\nabla\cdot\bm{u} < 0$: compression.

\paragraph{Deviatoric (shear) part --- traceless.}
The traceless part of $\varepsilon$ describes shape change without
volume change:
\begin{equation}
  e_{\alpha\beta} = \varepsilon_{\alpha\beta}
  - \tfrac{1}{3}(\nabla\cdot\bm{u})\,\delta_{\alpha\beta}.
  \label{eq:deviatoric}
\end{equation}
This is the \emph{shear strain}: it measures how the angles between
material lines change.

\begin{tcolorbox}[colback=yellow!6!white, colframe=orange!70!black,
  fonttitle=\bfseries\small, title={Notation: Three Roles of $e_{\alpha\beta}$},
  breakable]
The symbol $e_{\alpha\beta}$ carries three related but distinct meanings in
this section.
\begin{center}
\renewcommand{\arraystretch}{1.3}
\begin{tabular}{@{}lll@{}}
\toprule
Symbol & Object & Equation \\
\midrule
$\varepsilon_{\alpha\beta}$ & full strain tensor (symmetric) & (\ref{eq:strain}) \\
$e_{\alpha\beta} = \varepsilon_{\alpha\beta} - \tfrac{1}{3}(\nabla\!\cdot\!\bm{u})\delta_{\alpha\beta}$
  & deviatoric (traceless) strain & (\ref{eq:deviatoric}) \\
$\dot{e}_{\alpha\beta} \equiv \partial e_{\alpha\beta}/\partial t$
  & deviatoric strain \emph{rate} & (\ref{eq:viscoelastic}) \\
\bottomrule
\end{tabular}
\renewcommand{\arraystretch}{1}
\end{center}
In the VK functional (Eq.~\ref{eq:sigma_vk}) the relevant quantity is the
deviatoric strain \emph{rate} $\dot{e}_{\alpha\beta}$ (time derivative of the
traceless part), sometimes written $e^{\alpha\beta}$ in index-raising notation.
When reading the VK equation, $e^{\alpha\beta}$ always means the strain rate,
not the strain itself.
\end{tcolorbox}

\begin{derivbox}
\textbf{Concrete example: simple shear.}
Suppose the displacement is purely in the $x$-direction and varies
only with $y$: $\bm{u} = (u_x(y), 0, 0)$.  Then:
\[
  \varepsilon_{xy} = \varepsilon_{yx} = \frac{1}{2}\frac{\partial u_x}{\partial y},
  \quad \text{all other components zero.}
\]
The material is being sheared: layers at different $y$ slide past each
other.  A square element deforms into a parallelogram.  No volume change
($\nabla\cdot\bm{u}=0$), only angular distortion.
\end{derivbox}

In the VK functional, the relevant quantity is the \emph{strain rate
tensor} --- the time derivative of the strain:
\begin{equation}
  \dot{e}_{\alpha\beta}(\bm{r},t)
  = \frac{\partial e_{\alpha\beta}}{\partial t}
  = \frac{1}{2}\left(\frac{\partial v_\alpha}{\partial r_\beta}
                   + \frac{\partial v_\beta}{\partial r_\alpha}\right)
  - \frac{1}{3}(\nabla\cdot\bm{v})\,\delta_{\alpha\beta},
  \label{eq:strain_rate}
\end{equation}
where $\bm{v} = \partial\bm{u}/\partial t$ is the fluid velocity.

\subsubsection{The Stress Tensor: Internal Forces}

When a material is deformed, internal restoring forces develop.
The \emph{stress tensor} $\sigma_{\alpha\beta}$ describes these forces:
$\sigma_{\alpha\beta}$ is the force per unit area in the $\alpha$-direction
acting on a surface whose outward normal points in the $\beta$-direction.

\begin{itemize}
  \item \textbf{Diagonal elements} $\sigma_{\alpha\alpha}$: normal stresses
    (tension or compression along $\alpha$).
  \item \textbf{Off-diagonal elements} $\sigma_{\alpha\beta}$ ($\alpha\neq\beta$):
    shear stresses (forces parallel to the surface).
  \item \textbf{Symmetry}: $\sigma_{\alpha\beta} = \sigma_{\beta\alpha}$
    (required so that fluid elements have no net torque).
\end{itemize}

The equation of motion for a fluid element is Newton's second law
for a continuum:
\begin{equation}
  \rho_m \frac{\partial v_\alpha}{\partial t}
  = \sum_\beta\frac{\partial\sigma_{\alpha\beta}}{\partial r_\beta} + f_\alpha,
  \label{eq:newton_fluid}
\end{equation}
where $\rho_m$ is the mass density and $f_\alpha$ is any external body
force per unit volume.  For the electron gas, this becomes the force
balance between the XC stress tensor and the XC electric field
(Eq.~\ref{eq:sigma_xc}).

\subsubsection{Constitutive Relations: How Stress Relates to Strain}
\label{sec:constitutive}

A \emph{constitutive relation} specifies how the stress $\sigma_{\alpha\beta}$
depends on the deformation.  Different materials obey different constitutive
relations:

\paragraph{1. Elastic solid (Hooke's law).}
Stress is proportional to the instantaneous strain:
\begin{equation}
  \sigma_{\alpha\beta} = G_0\,e_{\alpha\beta}
  + K_0\,(\nabla\cdot\bm{u})\,\delta_{\alpha\beta},
  \label{eq:hookes_law}
\end{equation}
where $G_0$ is the shear modulus and $K_0$ is the bulk modulus.
No memory: the stress depends only on the current deformation.
If the force is removed, the material springs back instantly (rubber
on short timescales).

\paragraph{2. Viscous Newtonian fluid (Newton's law of viscosity).}
Stress is proportional to the instantaneous strain \emph{rate}:
\begin{equation}
  \sigma_{\alpha\beta}^{\rm visc}
  = \eta\,\dot{e}_{\alpha\beta}
  + \zeta\,(\nabla\cdot\bm{v})\,\delta_{\alpha\beta},
  \label{eq:newtons_viscosity}
\end{equation}
where $\eta$ is the \emph{shear viscosity} (resistance to sliding)
and $\zeta$ is the \emph{bulk viscosity} (resistance to compression).
Again no memory: only the current rate of deformation matters.
Water is well-described by Eq.~(\ref{eq:newtons_viscosity}) at low flow speeds.

\paragraph{3. Viscoelastic material: memory in the constitutive relation.}
A viscoelastic material has both elastic and viscous character.
The general constitutive relation replaces the instantaneous $G_0$
by a \emph{time-convolution with a relaxation kernel} $G(\tau)$:
\begin{equation}
  \sigma_{\alpha\beta}(t)
  = \int_0^\infty G(\tau)\,\dot{e}_{\alpha\beta}(t-\tau)\,\dd\tau
  = \int_{-\infty}^t G(t-t')\,\dot{e}_{\alpha\beta}(t')\,\dd t'.
  \label{eq:viscoelastic}
\end{equation}
The \emph{relaxation modulus} $G(\tau)$ is the kernel.  It tells you
how strongly the shear strain rate from $\tau$ seconds ago still
influences the current stress.

\begin{derivbox}
\textbf{Recovering elastic and viscous limits from the general form.}

\textbf{Elastic limit}: $G(\tau) = G_0$ (constant, infinite memory).
The material ``remembers'' its entire deformation history with equal
weight. Substituting into Eq.~(\ref{eq:viscoelastic}):
\[
  \sigma(t) = G_0\int_0^\infty \dot{e}(t-\tau)\,\dd\tau
  = G_0\bigl[e(t) - e(-\infty)\bigr]
  = G_0\,e(t),
\]
where the last step uses $e(-\infty)=0$ (the solid starts undeformed).
This is stress proportional to strain --- Hooke's law. $\checkmark$

\textbf{Viscous limit}: $G(\tau) = \eta\,\delta(\tau)$ (instantaneous
memory, no history). Substituting:
\[
  \sigma(t) = \int_0^\infty \eta\,\delta(\tau)\,\dot{e}(t-\tau)\,\dd\tau
  = \eta\,\dot{e}(t).
\]
Stress proportional to strain \emph{rate} --- Newton's law of viscosity.
$\checkmark$

\textbf{General viscoelastic}: $G(\tau)$ is a smoothly decaying
function, e.g.\ $G(\tau) = G_\infty + (G_0-G_\infty)e^{-\tau/\tau_{\rm rel}}$.
\begin{itemize}[noitemsep]
  \item For $\omega\tau_{\rm rel}\gg 1$ (fast deformation): the material
    has no time to relax --- behaves like an elastic solid.
  \item For $\omega\tau_{\rm rel}\ll 1$ (slow deformation): the material
    relaxes before the next cycle --- behaves like a viscous fluid.
\end{itemize}
Silly putty, polymer melts, and --- crucially --- the electron gas all
behave this way.
\end{derivbox}

\subsubsection{Summary: Mechanical Analogies for the VK Functional}

The Vignale-Kohn XC stress tensor (Eq.~\ref{eq:sigma_vk}) is
Eq.~(\ref{eq:viscoelastic}) applied to the electron gas:
\begin{itemize}[noitemsep]
  \item $e_{\alpha\beta}(t')$ $\to$ strain rate of the electron fluid
    at past time $t'$, computed from the current density $\bm{J}$
    (Eq.~\ref{eq:current_displacement}).
  \item $G(\tau)$ $\to$ $\eta_{\rm xc}(n_0,\tau)$ and $\zeta_{\rm xc}(n_0,\tau)$,
    the XC shear and bulk viscosity kernels of the homogeneous electron gas.
  \item $\sigma_{\alpha\beta}(t)$ $\to$ $\sigma_{\rm xc}^{\alpha\beta}(\bm{r},t)$,
    the XC stress that drives the correction to the KS vector potential.
\end{itemize}
The kernels $\eta_{\rm xc}$ and $\zeta_{\rm xc}$ are material properties
of the electron gas --- analogous to the viscosity of water --- that
must be computed from quantum Monte Carlo for the homogeneous electron gas
and then applied locally via the LDA of CDFT.

% -----------------------------------------------------------------------
\subsection{The Electron Fluid as a Viscoelastic Medium}

Vignale and Kohn showed that for a slowly varying perturbation (the
``gradient'' approximation in space), the XC contribution can be expressed
as a \emph{stress tensor} $\overleftrightarrow{\sigma}_{\rm xc}$, relating
to the XC vector potential via:
\begin{equation}
  n_0(\bm{r})\,\bm{E}_{\rm xc}(\bm{r},t) = \nabla\cdot\overleftrightarrow{\sigma}_{\rm xc}(\bm{r},t),
  \label{eq:sigma_xc}
\end{equation}
where $\bm{E}_{\rm xc} = -\partial\bm{A}_{\rm xc}/\partial t$.

\paragraph{Connection to your car analogy.}
This is the point where the analogy becomes exact.  The stress tensor has
the form of a viscoelastic constitutive relation:
\begin{equation}
  \sigma_{\rm xc}^{\alpha\beta}(\bm{r},t)
  = \int_{-\infty}^{t}
    \underbrace{\eta_{\rm xc}(n_0,\,t-t')}_{\text{shear viscosity kernel}}
    \underbrace{e^{\alpha\beta}(\bm{r},t')}_{\text{strain rate}}
    \,\dd t'
  + \int_{-\infty}^{t}
    \underbrace{\zeta_{\rm xc}(n_0,\,t-t')}_{\text{bulk viscosity kernel}}
    \underbrace{(\nabla\cdot\bm{u})\,\delta^{\alpha\beta}}_{\text{compression}}
    \,\dd t',
  \label{eq:sigma_vk}
\end{equation}
where:
\begin{itemize}
  \item $\bm{u}(\bm{r},t)$ is the \textbf{displacement field} of the electron
    fluid, related to the current by
    $\bm{J}(\bm{r},t) = n_0(\bm{r})\,\partial\bm{u}/\partial t$;
  \item $e^{\alpha\beta} = \tfrac{1}{2}(\partial_\alpha u_\beta +
    \partial_\beta u_\alpha) - \tfrac{1}{3}\delta_{\alpha\beta}\nabla\cdot\bm{u}$
    is the \textbf{symmetrized strain rate tensor} (shear deformation);
  \item $\eta_{\rm xc}(n_0,\tau)$ and $\zeta_{\rm xc}(n_0,\tau)$ are the
    \textbf{shear and bulk viscosity kernels} of the homogeneous electron gas
    at density $n_0$, computable from quantum Monte Carlo.
\end{itemize}

\begin{insight}
The strain rate $e^{\alpha\beta}$ encodes how fast the electron fluid is
being sheared at time $t'$: it is exactly your car analogy.  The XC
response today (stress $\sigma_{\rm xc}$ at time $t$) depends on the
full history of the deformation $e^{\alpha\beta}(\bm{r},t')$ for all
past times $t'\leq t$, weighted by the viscosity kernel.  A static
electron cloud produces zero shear and zero VK correction; a rapidly
oscillating current produces a large strain-rate history and a large
VK correction.
\end{insight}

\subsection{Viscosity Kernels: Explicit Forms and Limits}

The kernels are related to the frequency-dependent XC moduli of the
homogeneous electron gas:
\begin{align}
  \eta_{\rm xc}(n_0,\omega)
  &= -\frac{n_0^2}{\omega}\operatorname{Im}\fxc^T(n_0,\omega),
  \label{eq:eta_xc}\\
  \zeta_{\rm xc}(n_0,\omega)
  &= -\frac{n_0^2}{\omega}\left[\operatorname{Im}\fxc^L(n_0,\omega)
     - \tfrac{4}{3}\operatorname{Im}\fxc^T(n_0,\omega)\right],
  \label{eq:zeta_xc}
\end{align}
where $\fxc^T$ and $\fxc^L$ are the \textbf{transverse} and
\textbf{longitudinal} XC kernels of the uniform electron gas.
These labels refer to the direction of the density fluctuation relative
to its wavevector $\bm{q}$:
\begin{itemize}[noitemsep]
  \item \textbf{Longitudinal} ($\fxc^L$): the density fluctuation is
    \emph{parallel} to $\bm{q}$ --- a compressional (sound-wave) mode.
    This is the only component that survives in ordinary TDDFT based
    on the density $\rho$ alone, because $\rho$ is a scalar (longitudinal
    field).
  \item \textbf{Transverse} ($\fxc^T$): the density fluctuation is
    \emph{perpendicular} to $\bm{q}$ --- a shear mode (like vorticity
    in a fluid). Transverse response requires the vector current density
    $\bm{J}$ as the basic variable, which is why current-DFT (VK) is
    needed: ordinary TDDFT has $\fxc^T = 0$ by construction.
\end{itemize}
In the $q\to 0$ (long-wavelength) limit relevant to optical response,
both $\fxc^L$ and $\fxc^T$ approach finite values that are the XC
compressional and shear moduli of the electron gas.

The known limiting behaviors are:
\begin{align}
  \lim_{\omega\to 0}\operatorname{Im}\fxc^T(\omega) &\propto \omega
  &&\text{(Drude-like, weak low-frequency dissipation)},\\
  \lim_{\omega\to\infty}\operatorname{Im}\fxc^T(\omega) &\propto \omega^{-3/2}
  &&\text{(high-frequency algebraic decay)}.
\end{align}
In the time domain, $\eta_{\rm xc}(\tau)$ is a causal function that
decays over a characteristic time $\sim 1/\omega_p$ (the inverse plasma
frequency), meaning the memory has a finite range.  This is the basis for
the practical truncation of the history integral in numerical implementations.

% ============================================================================
\section{Functional Forms and Implementation of Memory Functionals}
\label{sec:implementation}
% ============================================================================

\subsection{The Simplest Memory Extension: Linear Convolution on Top of LDA}

The minimal memory extension retains the scalar $\Vxc$ form and adds a
linear convolution:
\begin{equation}
  \Vxc(\bm{r},t) = \Vxc^{\rm LDA}[\rho(\bm{r},t)]
  + \int_0^t K_{\rm xc}(n_0(\bm{r}),\,t-t')\,\delta\rho(\bm{r},t')\,\dd t',
  \label{eq:memory_lda}
\end{equation}
where $\delta\rho = \rho(t') - \rho_{\rm gnd}$ and $K_{\rm xc}(n_0,\tau)$
is the Fourier transform of $[\fxc(n_0,\omega) - \fxc(n_0,0)]$ ---
the frequency-dependent correction relative to the static (adiabatic) value.

This form:
\begin{itemize}[noitemsep]
  \item Reduces to adiabatic LDA when $K_{\rm xc}=0$ (or $\omega\to 0$).
  \item Captures shifts in excitation energies due to frequency-dependent
    screening.
  \item Does not capture full VK dissipation (which requires the vector
    potential $\bm{A}_{\rm xc}$).
\end{itemize}

\paragraph{Implementation cost.}
\begin{enumerate}[noitemsep]
  \item Store $\delta\rho(\bm{r},t')$ for $t' \in [t-\tau_{\rm max}, t]$:
    array of size $N_{\rm grid}\times N_{\rm history}$.
  \item At each step: convolve with $K_{\rm xc}(n_0(\bm{r}),\tau)$:
    $O(N_{\rm grid}\times N_{\rm history})$ operations.
  \item Truncate history at $\tau_{\rm max}\sim 10/\omega_p$
    (kernel is exponentially small beyond this).
  \item Memory overhead: $\sim N_{\rm grid}\times 100$ doubles (modest).
\end{enumerate}

\subsection{The Vignale-Kohn Functional: Full Implementation Requirements}

The full VK functional requires computing $\bm{A}_{\rm xc}$ from
the stress tensor via Eq.~(\ref{eq:sigma_xc}):
\begin{equation}
  -n_0\,\frac{\partial\bm{A}_{\rm xc}}{\partial t}
  = \nabla\cdot\overleftrightarrow{\sigma}_{\rm xc}.
  \label{eq:Axc_from_sigma}
\end{equation}
Integrating in time gives $\bm{A}_{\rm xc}(\bm{r},t)$, which then
enters the KS Hamiltonian as an additional velocity-gauge coupling:
\begin{equation}
  H_{\rm KS}(t) \to H_{\rm KS} + [\bm{A}_{\rm ext}(t)
  + \bm{A}_{\rm xc}(\bm{r},t)]\cdot\hat{\bm{p}}.
  \label{eq:H_vk}
\end{equation}

The matrix element of the XC correction in the KS basis is:
\begin{equation}
  \mel{\varphi_i}{\bm{A}_{\rm xc}\cdot\hat{\bm{p}}}{\varphi_j}
  = \int\!\bm{A}_{\rm xc}(\bm{r},t)\cdot
    \varphi_i^*(\bm{r})(-\ii\nabla)\varphi_j(\bm{r})\,\dd^3r.
  \label{eq:vk_matrix}
\end{equation}
This has the same structure as the external-field matrix element
$P^\alpha_{ij} = \ii v_{\rm cell}\mel{\varphi_i}{\partial_\alpha}{\varphi_j}$
already computed in the RMG velocity-gauge framework.  The VK correction
enters the same code path --- it is an additional contribution to the
effective $\bm{A}(t)$ at each timestep.

\begin{table}[ht]
\centering
\caption{Comparison of XC approximations for RT-TDDFT.}
\begin{tabular}{@{}lllll@{}}
\toprule
Functional & Memory & Dissipation & Doubles & Cost \\
\midrule
Adiabatic LDA/GGA (current) & No & No & No & $O(N^3)$ \\
Scalar memory [Eq.~(\ref{eq:memory_lda})] & Partial & No & Partial & $+O(N_h)$ \\
Vignale-Kohn (VK) & Yes (vector) & Yes & Partial & $+O(N_h)$ + Poisson \\
Exact TDDFT & Full & Full & Full & Exponential \\
\bottomrule
\end{tabular}
\end{table}

\subsection{Causality and Kramers-Kronig Relations}

Any physically correct memory kernel must satisfy the
\emph{Kramers-Kronig (KK) relations}.
These are not independent constraints imposed from outside --- they
follow automatically from causality via \textbf{complex analysis}.

\paragraph{Origin from causality and analyticity.}
A causal kernel $\fxc(\tau)=0$ for $\tau<0$ has a one-sided Fourier
transform (Eq.~\ref{eq:fxc_fourier}).
The key consequence is that $\fxc(\omega)$, extended to complex $\omega$,
is \emph{analytic in the upper half-plane} $\operatorname{Im}\omega > 0$:
no poles or branch cuts there, because the $e^{i\omega\tau}$ factor
is exponentially damped when $\operatorname{Im}\omega>0$ and $\tau>0$.

Applying the \textbf{Cauchy integral theorem} to a contour that runs
along the real axis and closes in the upper half-plane (where $\fxc\to 0$
at large $|\omega|$), one obtains the \textbf{dispersion relation}:
\begin{equation}
  \fxc(\omega) = \frac{1}{\pi\ii}\,\mathcal{P}\!\int_{-\infty}^\infty
  \frac{\fxc(\omega')}{\omega'-\omega}\,\dd\omega',
\end{equation}
where $\mathcal{P}$ denotes the Cauchy principal value.
Separating real and imaginary parts yields the Kramers-Kronig relations
($K_{\rm xc}(\tau) = 0$ for $\tau < 0$):
\begin{align}
  \operatorname{Re}\fxc(\omega)
  &= \fxc(0) + \frac{2}{\pi}\,\mathcal{P}\!\int_0^\infty
     \frac{\omega'\,\operatorname{Im}\fxc(\omega')}{{\omega'}^2 - \omega^2}\,\dd\omega',
  \label{eq:kk1}\\
  \operatorname{Im}\fxc(\omega)
  &= -\frac{2\omega}{\pi}\,\mathcal{P}\!\int_0^\infty
     \frac{\operatorname{Re}\fxc(\omega') - \fxc(0)}{{\omega'}^2-\omega^2}\,\dd\omega'.
  \label{eq:kk2}
\end{align}
The adiabatic approximation satisfies the $\omega\to 0$ limit of
Eq.~(\ref{eq:kk1}) trivially --- it is just the static kernel.
But it violates Eq.~(\ref{eq:kk2}) because $\operatorname{Im}\fxc^{\rm ATDDFT} = 0$
for all $\omega > 0$, while the right-hand side of Eq.~(\ref{eq:kk2})
is generally nonzero.  This is another way of seeing that the adiabatic
approximation is inconsistent with a causal memory.

% ============================================================================
\section{Summary}
\label{sec:summary}
% ============================================================================

The four questions addressed in this tutorial are resolved as follows:

\paragraph{1. What does the adiabatic approximation discard?}
It replaces the history-dependent XC kernel $\fxc(\omega)$ --- a
function with poles and a nonzero imaginary part --- by a frequency-independent
constant.  This is equivalent to assuming that the XC response is
instantaneous (zero memory time).  Time-derivative corrections improve
this approximation perturbatively but cannot reproduce poles.

\paragraph{2. Why is there no dissipation in adiabatic TDDFT?}
The fundamental reason is that the KS Hamiltonian is Hermitian and the
Magnus propagator is unitary, conserving the eigenvalues of $P(t)$
(occupation numbers) exactly.  Population relaxation ($T_1$) requires
a bath and is absent regardless of $\fxc$.  Pure dephasing ($T_2$)
requires $\operatorname{Im}\fxc(\omega)\neq 0$, which the adiabatic
approximation sets to zero.  Both effects are correctly attributed to
unitarity and the absence of memory, not just one or the other.

\paragraph{3. Why are double excitations missing?}
Pure doubly excited states are dipole-dark: the matrix element
$\mel{EE}{\hat{x}}{GG} = 0$ exactly (Eq.~\ref{eq:double_forbidden}).
In the interacting system, Coulomb mixing creates states with mixed
single+double character that are dipole-accessible.  Adiabatic TDDFT
places these states at wrong energies because the Casida basis
spans only single KS excitations.  Recovering the correct energy
requires a pole in $\fxc$ at the double-excitation frequency
(Eq.~\ref{eq:maitra_pole}).

\paragraph{4. Why is $(\rho,\bm{J})$ needed for periodic systems?}
$\rho(\bm{r},t)$ is gauge-invariant (Eq.~\ref{eq:rho_invariant}) and
cannot distinguish the presence or absence of a macroscopic vector
potential $\bm{A}(t)$ in a crystal.  $\bm{J}(\bm{r},t)$ is also
gauge-invariant (Eq.~\ref{eq:J_invariant}) and carries the missing
information.  The Vignale-Kohn theorem establishes $(\rho,\bm{J})$
as the correct basic variable pair.  The resulting XC functional has a
viscoelastic memory structure (Eq.~\ref{eq:sigma_vk}) that directly
formalizes the intuition that the electron system's response today
depends on the full history of its deformation.

% ============================================================================
\bibliographystyle{unsrt}
\bibliography{rmg_tddft_theory}
% ============================================================================

\end{document}
